% JCIM v3.0 development manuscript
% Revised from the supplied monolithic draft on 2026-08-21.
% The integrated Supporting Information begins after the bibliography.

\documentclass[journal=jcisd8,manuscript=article]{achemso}

\usepackage[T1]{fontenc}
\usepackage{amsmath,amssymb}
\usepackage{booktabs}
\usepackage{multirow}
\usepackage{array}
\usepackage{makecell}
\usepackage{graphicx}
\usepackage{xcolor}
\usepackage{tabularx}
\usepackage{longtable}
\usepackage{pdflscape}
\usepackage{float}
\usepackage{subcaption}
\usepackage{siunitx}

\graphicspath{{../figures/}{figures/}}
\newcommand{\TPSA}{TPSA}
\newcommand{\MolLogP}{MolLogP}
\newcommand{\ConLoss}{ConLoss}
\newcommand{\pendingresult}{\phantom{0.000 $\pm$ 0.000}}
\newif\ifdevelopmentdraft
\developmentdrafttrue

\author{Ziqi\,Wang}
\email{Ziqi.Wang21@student.xjtlu.edu.cn}
\affiliation{Department of Chemistry and Materials Science, School of Science, Xi'an Jiaotong-Liverpool University, Suzhou 215123, P. R. China}
\alsoaffiliation{Department of Chemistry, University of Liverpool, Liverpool L69 7ZD, UK}

\author{Minhao\,Shi}
\affiliation{Department of Chemistry and Materials Science, School of Science, Xi'an Jiaotong-Liverpool University, Suzhou 215123, P. R. China}
\alsoaffiliation{Institute of Infection, Veterinary and Ecological Sciences, University of Liverpool, Liverpool L69 3BX, UK}

\author{Pengrui\,Sun}
\affiliation{Department of Chemistry and Materials Science, School of Science, Xi'an Jiaotong-Liverpool University, Suzhou 215123, P. R. China}
\alsoaffiliation{Department of Chemistry, University of Liverpool, Liverpool L69 7ZD, UK}

\author{Zhenghao\,Wu}
\affiliation{Department of Chemistry and Materials Science, School of Science, Xi'an Jiaotong-Liverpool University, Suzhou 215123, P. R. China}
\alsoaffiliation{Department of Chemistry, University of Liverpool, Liverpool L69 7ZD, UK}

\author{Xuan\,Xue}
\email{Xuan.Xue@xjtlu.edu.cn}
\affiliation{Department of Chemistry and Materials Science, School of Science, Xi'an Jiaotong-Liverpool University, Suzhou 215123, P. R. China}
\alsoaffiliation{Department of Chemistry, University of Liverpool, Liverpool L69 7ZD, UK}

\title{Descriptor-Guided Chemical Space Partitioning for Molecular Representation Learning}

\abbreviations{ADME, ConLoss, FLOPs, GNN, GPU, KAN, MLM, MLP, MoE, MPP, PCA, QSAR, QSPR, ROC-AUC, SMILES, TPSA, t-SNE}
\keywords{molecular pretraining, molecular property prediction, mixture of experts, physicochemical descriptors, MolLogP, topological polar surface area, Kolmogorov-Arnold networks, contrastive learning}

\begin{document}

\begin{tocentry}
\includegraphics[width=8.0cm]{toc_graphic_vector.pdf}
\end{tocentry}


\begin{abstract}
Molecular descriptors are routinely used as predictive features or auxiliary targets, but they rarely determine which parameters process a molecule during representation learning. We investigate descriptor-guided chemical space partitioning in which fixed intervals of Wildman-Crippen MolLogP or topological polar surface area (TPSA) select one final-layer computation branch after shared graph and SMILES pretraining, with branch-wise multilayer perceptron or Kolmogorov-Arnold network functions and descriptor-aware contrastive regularization. Historical screening selected final-layer conditioning and descriptor-specific mixed branch allocations, while the legacy eight-layer-aligned development snapshot gives macro area under the receiver operating characteristic curve (ROC-AUC) values of 0.7985 for the matched local Base, 0.7983 for TPSA-conditioned multilayer perceptron branches with contrastive regularization, 0.7810 for the complete TPSA model, and 0.7823 for the complete MolLogP model. Contrastive regularization substantially increases descriptor-region separation in the molecular embeddings, indicating that a physicochemical coordinate can organize representation geometry even when transfer to downstream property prediction is heterogeneous.
\end{abstract}

\ifdevelopmentdraft
\begin{center}
\fcolorbox{orange!70!black}{orange!7}{%
\parbox{0.94\textwidth}{\small\textbf{Development-status note.}
The numerical results in this draft are preserved from the legacy downstream
protocol. That protocol produced all 440 planned fold files, but validation used
shuffled batches with \texttt{drop\_last=True} and test evaluation used
\texttt{drop\_last=True}. File completeness therefore does not establish full
split coverage. Historical layer and KAN screens are treated as exploratory
design studies. A corrected full-validation and full-test evaluation is outside
the present manuscript-revision task and is required before the benchmark
results are described as final.}}
\end{center}
\fi



\section{Introduction}
\label{sec:introduction}

Molecular property prediction (MPP) supports compound prioritization throughout
drug discovery and chemical risk assessment. Models are used to estimate
physicochemical properties, absorption, distribution, metabolism, excretion,
toxicity, and target-associated activity before every candidate can be tested in
an assay. This computational triage does not remove the need for experiment. It
helps allocate experimental capacity to a chemical space that is much larger
than any practical screening collection. Deep-learning models have identified
antibacterial candidates outside established structural series and have been
used to examine structural features associated with predicted activity
\cite{Stokes2020,Wong2024}. Their practical value depends on both discrimination
within the available data and transfer to scaffolds or property regimes that
are weakly represented during training \cite{Walters2021,Segal2025}. Benchmark
design, split construction, and uncertainty across repeated runs are therefore
part of the chemical question rather than secondary implementation details
\cite{Wu2018,Sultan2024}.

Quantitative structure-activity relationship (QSAR) and quantitative
structure-property relationship (QSPR) modelling provide the chemical basis for
this task. Classical QSAR describes an endpoint through structural and
physicochemical variables whose relevance depends on the biological or physical
process being modelled \cite{Hansch1964,Cherkasov2014,Tropsha2024}. Lipophilicity
affects partitioning between aqueous and nonpolar environments, desolvation,
nonspecific association, and access to hydrophobic sites. Polarity and
hydrogen-bonding capacity influence solvation, intermolecular recognition, and
passive transport. Molecular size, flexibility, aromaticity, formal charge,
shape, and electronic structure provide further constraints. These factors do
not operate independently, and no single descriptor is expected to explain a
general molecular endpoint. The QSAR perspective nevertheless establishes an
important modelling principle. The same local structure can contribute
differently when it occurs in a different whole-molecule physicochemical
context.

This context also makes chemical space heterogeneous in a way that scaffold
identity alone does not capture. Molecules sharing a ring system can span
substantial ranges of ionizable functionality, polar surface, and
lipophilicity after substitution. Conversely, molecules with different
scaffolds can occupy similar physicochemical regimes and encounter related
solvation or transport constraints. A structural encoder and a descriptor
panel therefore describe overlapping but nonidentical aspects of a molecular
collection. The former resolves how atoms are connected and arranged. The
latter locates the complete molecule on experimentally motivated property
axes. A global coordinate can consequently condition how a structural
embedding is transformed without replacing the structural representation.

Molecular machine learning approaches structure-property relations through
learned representations. Extended-connectivity fingerprints encode circular
atomic neighbourhoods and remain strong cheminformatics baselines
\cite{Rogers2010}. Graph neural networks (GNNs) update atom representations from
their bonded environments and can learn task-dependent features without an
explicit descriptor panel \cite{Gilmer2017,Yang2019}. Attention, graph
Transformers, and geometry-aware networks extend this representation to longer
range or three-dimensional interactions \cite{Atz2021,Luong2024Cheminformatics}.
Molecular pretraining uses larger unlabeled collections to fit an encoder before
property-specific labels are introduced. Masked graph reconstruction,
contextual prediction, and graph-level contrastive learning have produced
transferable representations across molecular benchmarks
\cite{Hu2020,Rong2020,Wang2022,Fang2022}. Multimodal approaches combine a
molecular graph with a simplified molecular input line entry system (SMILES)
sequence so that connectivity and sequence context provide complementary
training signals \cite{Ross2022MolFormer,Shen2024MoleSG}. These models capture
atom environments, bond patterns, scaffolds, and other structural regularities
that would be difficult to specify exhaustively as hand-engineered features.

Physicochemical descriptors remain common in this learned-representation
setting, although their computational role is usually limited. A descriptor can
be used directly in a classical model, concatenated with a learned embedding, or
predicted as an auxiliary pretraining target \cite{Shen2021MolMap,Li2023}.
Recent studies have used two-dimensional descriptors, quantum properties, and
other property targets to enrich molecular pretraining
\cite{Kim2024QuantumInformed,Li2024QMDescriptors,Chang2024SPMM}. These strategies
allow physicochemical information to influence a representation or a final
predictor. They generally do not use a descriptor to determine which
parameterized transformation processes the molecule. Two molecules with
different polarity or lipophilicity can therefore pass through the same
feed-forward function even if the latent representation contains information
correlated with those properties. Encoding a descriptor and using it to
organize computation are distinct modelling choices.

That distinction has a practical chemical interpretation. Feature
concatenation asks whether a predictor can use a numerical descriptor after
representation learning. Auxiliary prediction asks whether the latent space
contains enough information to reconstruct it. Descriptor-conditioned
computation asks whether molecules in different physicochemical regions
benefit from different parameterized transformations while retaining a shared
structural encoder. The descriptor then acts as an assignment variable with
known units and boundaries. Region occupancy can be measured in both
pretraining and application data, and the active branch can be recovered for
an individual molecule. These properties do not guarantee better prediction.
They make the relation between a global chemical coordinate and model
execution experimentally accessible.

We study the second choice through descriptor-guided chemical space
partitioning. Let \(\mathcal{M}\) denote the molecular domain and let a scalar
descriptor be a mapping \(d:\mathcal{M}\rightarrow\mathbb{R}\). For molecule
\(m_i\), the coordinate \(z_i=d(m_i)\) is assigned to one fixed interval in
\(\mathcal{P}_d=\{I_0,I_1,\ldots,I_{K-1}\}\). The assignment is
\begin{equation}
e_i=\phi_d(z_i),
\qquad
\phi_d(z_i)=k\ \text{when}\ z_i\in I_k.
\end{equation}
The resulting region index selects one computation branch. Shared graph and
SMILES layers continue to learn molecular structure, while the descriptor
controls the nonlinear function used at the conditioned layer. This formulation
is descriptor-agnostic in the limited technical sense that another scalar
descriptor could define \(\phi_d\). It does not imply that every descriptor,
interval scheme, or chemical domain will provide a useful partition.

Wildman-Crippen MolLogP and topological polar surface area (TPSA) provide two
chemically established instances. MolLogP is an atom-contribution estimate of
lipophilicity for the supplied two-dimensional molecular structure
\cite{Wildman1999}. Lipophilicity is relevant to aqueous and nonpolar
partitioning, solubility, membrane access, and nonspecific binding, which is why
calculated logP appears frequently in medicinal-chemistry property analysis
\cite{Lipinski1997,Bickerton2012}. MolLogP is not a pH-dependent distribution
coefficient and does not enumerate microspecies. TPSA is a fragment-based sum
of contributions from polar atoms and attached hydrogens \cite{Ertl2000}. It is
related to polarity, hydrogen-bonding interactions, desolvation, permeability,
and oral bioavailability \cite{Veber2002}. TPSA is not a direct electronic
descriptor such as a substituent constant, and it does not describe
conformer-dependent solvent exposure. The two descriptors are inexpensive,
deterministic, broadly available in RDKit, and interpretable across common
drug-like structures. They probe complementary lipophilicity-related and
polarity-related coordinates without being presented as a complete molecular
property basis.

Their joint use also provides a control on descriptor-specific interpretation.
MolLogP and TPSA correlate with several other molecular attributes, including
heteroatom content, hydrogen-bond counts, aromaticity, and size, yet they order
the same corpus differently. A branch selected by MolLogP therefore cannot be
assumed to correspond to the matching TPSA branch or to a single structural
category. Comparing the two axes tests whether an observed computation pattern
follows one chosen coordinate or persists across chemically distinct
partitions. Molecular weight, flexibility, ionization-aware descriptors, and
three-dimensional surface measures remain plausible future coordinates. Each
would require its own boundary design, occupancy audit, and downstream
evaluation.

The computation branches are implemented within a mixture of experts (MoE)
layer. In conventional MoE models, a learned router maps a hidden representation
to one or more modules \cite{Jacobs1991,Shazeer2017,Fedus2022}. Molecular MoE
methods have used learned topology, scaffolds, or substructure representations
to organize predictions \cite{Kim2023,Jiang2025ESIBMol}. Such routing can adapt
to the training objective, but the assignment itself is latent and can change
during optimization. Here, the descriptor value and interval boundaries fix
the assignment before training. Following MoE terminology, the branches are
implemented as expert modules. Their scientific interpretation is limited to
parameterized functions associated with descriptor-defined regions. A branch
does not represent a discrete chemical class. The main contribution is
therefore not a new generic MoE architecture. It is an experiment on whether a
named physicochemical coordinate can determine a molecular computation path.

Descriptor-based assignment leaves open where that path should diverge from
shared structural processing. Early graph Transformer layers encode atoms,
bonds, and local neighbourhoods that occur across the descriptor range. Applying
separate functions at every layer may divide this common processing before a
higher-level molecular representation has formed. We therefore compare final,
odd, and even layer placements. The comparison asks where descriptor-conditioned
computation should enter, rather than which generic MoE depth is universally
best.

A second question concerns the nonlinear function used in each region. Fixed
intervals differ in population size, scaffold composition, and correlated
molecular properties. Uniform replacement of every multilayer perceptron (MLP)
with a Kolmogorov-Arnold network (KAN) would assume that all regions benefit from
the same function family. KANs replace scalar edge weights with learnable
univariate functions, and KAN modules have recently been studied in molecular
graph models \cite{Liu2025,Li2025}. We screen one branch at a time against the
matching all-MLP partition. This branch-wise procedure tests regional function
choice under a fixed descriptor axis. It does not assume that KAN is generally
more expressive, more interpretable, or more efficient for molecular data.

Assignment also does not guarantee that the final molecular embedding retains
the descriptor partition. A molecule can traverse one branch while subsequent
optimization produces overlapping region distributions. We therefore add a
descriptor-aware contrastive loss (ConLoss) in which molecules assigned to the
same region form positive pairs \cite{Khosla2020}. ConLoss provides an explicit
representation-level objective. Its effect is evaluated separately from
property prediction because embedding organization and endpoint discrimination
are not interchangeable \cite{Zhang2024RepresentationAudit,Welsch2026}. The
two-dimensional t-distributed stochastic neighbor embedding (t-SNE) plots are
used only for visualization, while region-separation statistics are calculated
in principal-component space.

The experimental sequence follows these dependencies. We first characterize
the occupancy and chemistry of every fixed interval in the ZINC15 pretraining
corpus. We then examine where descriptor-conditioned computation is introduced
and whether individual regions favour MLP or KAN functions. The frozen
configurations are compared across eight scaffold-split classification
datasets. Conditional ablation separates the effect of descriptor-conditioned
branches, partial KAN, ConLoss, and their interaction. Embedding geometry,
stored capacity, active arithmetic, and observed runtime are reported as
separate outcomes.

The study makes three contributions to molecular representation analysis.
First, it formulates a molecular descriptor as a deterministic selector of
conditional computation rather than an appended feature or auxiliary target.
Second, it evaluates branch-specific nonlinear function selection within fixed
TPSA and MolLogP regions. Third, it measures whether descriptor-aware
regularization is reflected in molecular embedding organization and whether
that organization transfers to downstream prediction. The current evidence
supports traceable computation and strong control of descriptor-region geometry.
It does not support a claim of uniform predictive improvement, which sets the
boundary for the results and discussion that follow.

\clearpage

\section{Methods}
\label{sec:methods}

\subsection{Data Sources, Molecular Curation, and Evidence Status}

Pretraining uses the 250,000-molecule ZINC15 subset distributed with the
MoleSG workflow \cite{Sterling2015ZINC,Shen2024MoleSG}. The existing
preprocessing reads the supplied RDKit molecule objects, assigns stereochemical
tags from the stored structures, computes Gasteiger charges for atom features,
and generates a canonical SMILES string with \texttt{Chem.MolToSmiles}. No
pH-dependent microspecies generation, tautomer enumeration, or conformer-based
surface calculation is introduced for the descriptor assignment in this study.
MolLogP and TPSA are calculated from the same two-dimensional RDKit structure
that enters preprocessing. Descriptor values therefore refer to that explicit
representation rather than to an experimental logD or a conformer-dependent
solvent-accessible surface.

The downstream datasets are the eight MoleculeNet classification collections
listed in Table~\ref{tab:datasets} \cite{Wu2018}. Labels equal to \(-1\) are
retained as missing-value markers and are masked during endpoint evaluation.
No endpoint labels are generated or imputed. The current manuscript preserves
the result files produced by the historical downstream protocol. A separate
code audit found that the historical validation loader used shuffled batches
with \texttt{drop\_last=True}, while the test loader used deterministic order
with \texttt{drop\_last=True}. The legacy result namespace contains every
planned model, dataset, and fold file, but those files do not necessarily cover
every row in the corresponding validation or test split. The numerical tables
are consequently reported as a legacy development snapshot. No corrected
training or inference is introduced during this manuscript revision.

\subsection{Molecular Corpus and Backbone Pretraining}

The model extends the multimodal self-supervised MoleSG framework, which combines molecular graphs with SMILES strings \cite{Shen2024MoleSG}.

Pretraining uses 250,000 unlabeled ZINC15 molecules \cite{Sterling2015ZINC}. Each molecule is represented by a graph and a tokenized SMILES sequence. Atom feature vectors have length 115, and bond feature vectors have length 13. Nonoverlapping corruption is applied to the two modalities. Approximately one quarter of graph atoms are selected for masked atom reconstruction. SMILES tokens that do not correspond to graph-masked atoms are eligible for masked language modeling (MLM).

The graph decoder reconstructs masked atom targets with scaled cosine error (SCE). For normalized target \(\widehat{\mathbf y}_j\) and prediction \(\widehat{\mathbf p}_j\), the implementation uses
\begin{equation}
\mathcal{L}_{\mathrm{SCE}}
=
\frac{1}{N_{\mathrm{mask}}}
\sum_{j=1}^{N_{\mathrm{mask}}}
\left(
1-\widehat{\mathbf y}_j^{\mathsf T}\widehat{\mathbf p}_j
\right)^{\alpha},
\qquad
\alpha=1.
\end{equation}
The SMILES decoder uses token-level cross-entropy over a 700-token vocabulary. The graph encoder contains eight Transformer layers with hidden dimension 256 and eight attention heads. Each standard feed-forward network contains two dense transformations. The encoder uses ScaleNorm, an exponential distance kernel, and zero dropout. Unless otherwise stated, pretraining uses batch size 32, seed 42, 300 epochs, and Adam with a Transformer learning-rate schedule. The graph and SMILES encoders are connected through the original MoleSG multimodal module. The base pretraining objective is
\begin{equation}
\mathcal{L}_0
=
\mathcal{L}_{\mathrm{SCE}}
+
\mathcal{L}_{\mathrm{MLM}}.
\end{equation}

\subsection{Descriptor-Defined Regions}

Let \(d\) denote a scalar molecular descriptor. Molecule \(m_i\) receives descriptor value \(z_i=d(m_i)\) and region index \(e_i=\phi(z_i)\). The assignment function \(\phi\) is defined by fixed intervals and contains no learned parameters. Descriptor calculations use RDKit version 2022.09.5 \cite{RDKit2022}. TPSA is calculated with the Ertl fragment method, and MolLogP is calculated with the Wildman-Crippen atom-contribution model.

TPSA boundaries are 20, 40, 60, 80, 100, 120, and \(140\ \text{\AA}^2\). MolLogP boundaries are \(-1\), 0, 1, 2, 3, 4, and 5. Values below the first boundary enter region 0. Values at or above the final boundary enter region 7. Intermediate intervals are left-closed and right-open. These boundaries are fixed chemical intervals rather than empirical quantiles, so region occupancy is determined by the descriptor distribution of the pretraining corpus.

Pretraining reads one stored region index per ZINC15 source row. The source-row index is recovered from each preprocessing filename and is used to index the assignment array. For this manuscript, TPSA and MolLogP were recomputed for all 250,000 source rows and reassigned with the stated boundaries. Both stored arrays matched the recomputed indices without mismatches, and the preprocessing directory contained one indexed file for every source row. Equal-frequency bins used in the earlier descriptor-screening script were limited to descriptor ranking and were not used by the pretraining assignment file.

Descriptor-region profiles are calculated over the complete pretraining corpus. Molecular weight uses RDKit \texttt{MolWt}. Hydrogen-bond donor and acceptor counts use \texttt{NumH\-Donors} and \texttt{NumH\-Acceptors}. Rotatable bonds use \texttt{Num\-Rotatable\-Bonds}, and ring count uses \texttt{Calc\-Num\-Rings}. Aromatic atom fraction is the number of aromatic atoms divided by the number of heavy atoms. Heteroatom fraction is the number of nonhydrogen and noncarbon atoms divided by the number of heavy atoms. Molecular formal charge is the sum of atomic formal charges. Each region is reported with its molecule count, corpus fraction, arithmetic mean, and sample standard deviation. Heatmap values are z scores of the eight region means for each chemical variable. Regions containing fewer than 10 molecules are masked in the profile heatmap. Descriptor correlations use Pearson correlation on a deterministic 10\% sample obtained by retaining every tenth corpus entry.

\subsection{Descriptor-Conditioned Branches}

Let \(\mathbf{H}_i^{(\ell)}\) denote the node representation of molecule \(m_i\) at layer \(\ell\). At a descriptor-conditioned layer,
\begin{equation}
\widetilde{\mathbf{H}}_i^{(\ell)}
=
E_{e_i}^{(\ell)}\left(\mathbf{H}_i^{(\ell)}\right).
\end{equation}
One branch is evaluated for each molecule, and its region index is applied to every valid graph node of that molecule. An MLP branch contains two 256-to-256 linear transformations with Mish activation. All branches are stored in the model, while one branch is active for each molecule.

\begin{figure*}[t]
\centering
\includegraphics[width=\textwidth]{fig01_descriptor_partition_overview.pdf}
\caption{Descriptor-guided chemical space partitioning and conditional
computation. Neutral ibuprofen is used as a verified illustrative molecule
and is not presented as a sampled ZINC15 record. Its canonical SMILES is
\texttt{CC(C)Cc1ccc(C(C)C(=O)O)cc1}. RDKit 2022.09.5 gives TPSA
\(=37.30\ \text{\AA}^2\), which assigns zero-based region 1 and its KAN
branch, and MolLogP \(=3.073\), which assigns zero-based region 5 and its MLP
branch. MolLogP denotes the Wildman-Crippen value for the drawn neutral
structure rather than pH-dependent LogD. TPSA denotes topological polar
surface area rather than conformer-dependent exposed surface. The first seven
graph Transformer layers are shared and the eighth
layer contains the descriptor-conditioned branch. Property prediction and
descriptor-region organization are evaluated separately. The nonquantitative
visual background was generated with GPT Image2 and edited by the authors. The
molecular drawing, descriptor values, interval assignments, branch functions,
labels, and arrows were generated from RDKit and project code.}
\label{fig:workflow}
\end{figure*}

\subsection{Placement of Conditional Computation}

The exploratory architecture screen compares three sets of conditioned
Transformer layers. Final-layer placement replaces layer 8. Odd-layer
placement replaces layers 1, 3, 5, and 7. Even-layer placement replaces layers
2, 4, 6, and 8. All three configurations use the MolLogP partition and eight
MLP branches. The historical screen used downstream test summaries from the
legacy protocol. It is therefore used to document model development rather
than as an independent confirmatory test. Final-layer placement is retained as
the locked design in this revision. No new placement is selected from the
reported test results.

\subsection{Branch-Wise KAN Selection}

A KAN branch contains two KAN linear transformations. Each transformation combines a SiLU base mapping with a B-spline mapping. The grid size is 5 and the spline order is 3. For each descriptor partition, eight matched pretraining models are produced by replacing one branch \(k\) with KAN while retaining MLP in the remaining seven branches.

The historical branch screen uses the equally weighted mean ROC-AUC across
BBBP, Tox21, ToxCast, SIDER, ClinTox, BACE, HIV, and MUV. Five scaffold seeds
are available for each model and dataset combination except MolLogP branch 5,
which contains four folds and is retained as MLP. A branch was assigned KAN
when its legacy eight-dataset macro ROC-AUC was higher than the corresponding
all-MLP final-layer model. As with placement, these test-derived values are
reported as exploratory design evidence. They are not reused to claim an
independent test-set advantage.

This procedure assigns KAN to TPSA regions 0, 1, 2, 3, and 5. Regions 4, 6, and 7 retain MLP. For MolLogP, regions 1, 2, 3, and 4 use KAN, while regions 0, 5, 6, and 7 retain MLP. Each selected mixed model is then pretrained jointly from random initialization. Branch weights are not assembled from the individual screening checkpoints.

\subsection{Descriptor-Aware Contrastive Loss}

Graph node representations are mean-pooled over valid nodes and L2 normalized. For anchor molecule \(i\), molecules in the same descriptor region form the positive set. All other nonself molecules contribute to the contrastive denominator. Self-pairs are excluded. At most 4096 anchors are retained by the implemented loss. With temperature \(\tau=0.1\), the loss is
\begin{equation}
\mathcal{L}_{\mathrm{Con}}
=
-\frac{1}{|\mathcal{A}|}
\sum_{i\in\mathcal{A}}
\frac{1}{|\mathcal{P}(i)|}
\sum_{p\in\mathcal{P}(i)}
\log
\frac{
\exp\left(\mathbf{u}_i^{\mathsf{T}}\mathbf{u}_p/\tau\right)
}{
\sum_{a\neq i}
\exp\left(\mathbf{u}_i^{\mathsf{T}}\mathbf{u}_a/\tau\right)
}.
\end{equation}
Anchors without a positive molecule in the batch are excluded. The complete pretraining objective is
\begin{equation}
\mathcal{L}
=
\mathcal{L}_{\mathrm{SCE}}
+
\mathcal{L}_{\mathrm{MLM}}
+
0.1\mathcal{L}_{\mathrm{Con}}.
\end{equation}
Models without ConLoss set its coefficient to zero. The Supporting Information provides the full mathematical specification.

\subsection{Legacy Downstream Evaluation and ROC-AUC}

The result snapshot uses the eight classification datasets listed in
Table~\ref{tab:datasets}. Each dataset is divided into training, validation,
and test sets with a balanced scaffold split of approximately 0.8, 0.1, and
0.1. Splits are repeated with seeds 42, 43, 44, 45, and 46. Validation ROC-AUC
is used for early stopping with a patience of 20 epochs. The checkpoint from
the best observed validation epoch is evaluated once on the corresponding
test partition.

All models in the main snapshot use an eight-layer downstream graph encoder
with the same hidden dimension, attention heads, feed-forward depth,
ScaleNorm, distance kernel, and dropout as the pretraining backbone. Checkpoint
validation rejects missing or unexpected graph-encoder parameters. The earlier
placement and KAN screens used historical dataset-specific downstream settings
and remain separate from the eight-layer-aligned comparison.

Fold indices 1 through 5 correspond to seeds 42 through 46. In the historical
runs, training and validation use \texttt{shuffle} set to true. Test evaluation
uses \texttt{shuffle} set to false. All three loaders use \texttt{drop\_last}
set to true. Reported ROC-AUC values are recalculated from the
rows actually written to each legacy result CSV. Thus, 440 available fold files
mean that the planned experiment matrix is file-complete. They do not mean that
every split row was evaluated. Because only one historical best checkpoint is
retained per fold, the best epoch under complete validation coverage cannot be
reconstructed exactly by inference alone.

The corrected evaluation design retains incomplete final batches and uses
deterministic validation and test order. It also uses stable sample
identifiers, model selection based only on validation results, and one test
evaluation after the model definition is frozen.
Those corrected results are not yet available and are not inferred from the
legacy files.

For dataset \(d\), fold \(f\), and endpoint \(t\), labels equal to \(-1\) are masked. An endpoint is retained only when the observed test labels contain both classes. Its area under the receiver operating characteristic curve (ROC-AUC) is
\begin{equation}
A_{dft}
=
\operatorname{ROC\mbox{-}AUC}
\left(
\mathbf y_{dft},
\mathbf s_{dft}
\right).
\end{equation}
For a multitask dataset, the fold value gives equal weight to the valid endpoints
\begin{equation}
A_{df}
=
\frac{1}{|T_{df}|}
\sum_{t\in T_{df}}
A_{dft}.
\end{equation}
The dataset mean and sample standard deviation are
\begin{equation}
\overline{A}_d
=
\frac{1}{5}\sum_{f=1}^{5}A_{df},
\qquad
s_d
=
\sqrt{
\frac{1}{4}
\sum_{f=1}^{5}
\left(
A_{df}-\overline{A}_d
\right)^2
}.
\end{equation}
The eight-dataset macro ROC-AUC is the arithmetic mean of the eight dataset means. It is not weighted by the number of molecules or endpoints.

\begin{table*}[t]
\centering
\caption{Downstream classification datasets and the legacy local evaluation
protocol. Each scaffold split uses seeds 42, 43, 44, 45, and 46. File-complete
fold coverage does not imply full validation or test sample coverage because
the historical loaders used \texttt{drop\_last=True}.}
\label{tab:datasets}
\begin{tabular}{lrrrrrr}
\toprule
Dataset & Molecules & Tasks & Type & Split & Seeds & Metric \\
\midrule
BBBP & 2,035 & 1 & Classification & Scaffold & 5 & ROC-AUC \\
Tox21 & 7,821 & 12 & Classification & Scaffold & 5 & ROC-AUC \\
ToxCast & 8,577 & 617 & Classification & Scaffold & 5 & ROC-AUC \\
SIDER & 1,379 & 27 & Classification & Scaffold & 5 & ROC-AUC \\
ClinTox & 1,468 & 2 & Classification & Scaffold & 5 & ROC-AUC \\
BACE & 1,513 & 1 & Classification & Scaffold & 5 & ROC-AUC \\
HIV & 41,127 & 1 & Classification & Scaffold & 5 & ROC-AUC \\
MUV & 93,087 & 17 & Classification & Scaffold & 5 & ROC-AUC \\
\bottomrule
\end{tabular}
\end{table*}

\subsection{Paired Comparisons and Statistical Analysis}

Each descriptive result cell requires all five available legacy fold files.
For left model \(L\) and right model \(R\), the matched fold difference is
\begin{equation}
\delta_{df}
=
A_{df}^{(L)}
-
A_{df}^{(R)}.
\end{equation}
The dataset difference is the mean of its five fold differences. A dataset is counted as a win when this value is positive. The macro difference gives equal weight to the eight dataset differences.

Hierarchical paired bootstrap intervals use 10,000 iterations and random seed 42. Each iteration first samples eight datasets with replacement. For each selected dataset, five matched fold differences are sampled with replacement. Dataset means are calculated before the macro mean. The 2.5th and 97.5th percentiles define the 95\% confidence interval. The Supporting Information also reports two-sided Wilcoxon signed-rank tests over the eight dataset-level differences with \texttt{zero\_method=wilcox}. Holm correction is applied over the planned comparisons.

The ablation uses three component labels. \(I_1\) denotes descriptor-conditioned final-layer branches with MLP functions. \(I_2\) denotes branch-wise partial KAN. \(I_3\) denotes ConLoss. Partial KAN and ConLoss require \(I_1\), so the component analysis is conditional rather than a complete independent factorial design.
\begin{align}
\Delta I_1 &= I_1-\mathrm{Base},\\
\Delta I_2 &= (I_1+I_2)-I_1,\\
\Delta I_3 &= (I_1+I_3)-I_1,\\
\Delta_{\mathrm{interaction}}
&=
\mathrm{Full}-(I_1+I_2)-(I_1+I_3)+I_1.
\end{align}

\subsection{Computational Analysis}

Stored parameters include shared tensors and every stored branch tensor. Active parameters describe tensors used by one forward pass. For mixed branch models, the expected active count and arithmetic cost are weighted by descriptor-region occupancy in the ZINC15 pretraining corpus. Analytical feed-forward multiply-accumulates are calculated from the implemented layer dimensions. One multiply-accumulate is counted as two floating-point operations (FLOPs). KAN estimates include the SiLU base mapping and spline coefficient multiplications. They are lower bounds because B-spline basis construction and dispatch overhead are not represented as dense matrix multiplication. Learned-router controls include router projections and active top-\(k\) MLP branches. Router softmax and dispatch overhead are excluded.

Observed pretraining time is recovered from the in-process epoch timers. A complete epoch timer measures wall time within a running Slurm segment. The first epoch of every segment is excluded because it includes startup and data-loading initialization. Within each segment, durations above \(Q_3+3\,\mathrm{IQR}\) are also excluded. Queue waits, pauses, and intervals between continuation jobs are outside the epoch timer and do not enter active-time summaries. The stored \texttt{train loss} value is the composite loss from the final mini-batch of each epoch rather than an epoch average. It is reported as logged terminal-batch composite loss. Completed graphics processing unit (GPU) hours and unsuccessful or cancelled GPU hours are tabulated separately. Finetuning runtime records describe Slurm execution history and are not treated as a controlled architecture speed benchmark.

\subsection{Embedding Analysis and t-SNE}

The representation analysis is performed separately for every dataset, descriptor partition, and fold. Coordinate spaces from different folds are not combined. SMILES present in Base, I1, I1+I2, I1+I3, and Full define the common molecular set. For molecule \(i\), the RDKit atom count \(n_i\) identifies the valid rows of the padded node tensor. These rows are mean-pooled to obtain one molecular embedding.

Each model embedding matrix is independently reduced by principal component analysis (PCA) to 50 components or the largest dimension allowed by the sample count and feature dimension. PCA features are L2 normalized. Silhouette score uses Euclidean distance and compares within-region cohesion with separation from the nearest alternative region. Davies-Bouldin index is the mean, over regions, of the largest ratio between paired within-region dispersion and centroid separation. Lower values indicate better separation. Regions containing fewer than two molecules are excluded from these calculations.

Within-region distance is the mean nonself Euclidean pairwise distance among molecules assigned to the same retained region. Between-region distance is the mean distance among pairs assigned to different retained regions. Their ratio is between-region distance divided by within-region distance. These quantities and the clustering indices are calculated in PCA space. Five-fold summaries report the arithmetic mean and sample standard deviation only when all five folds are available.

The t-distributed stochastic neighbor embedding (t-SNE) visualization uses fold 1, two dimensions, perplexity 30, PCA initialization, learning rate \texttt{auto}, 2000 iterations, and random seed 42. HIV and MUV use deterministic descriptor-region-stratified samples of at most 2000 common molecules. Other datasets use all common molecules. The two-dimensional coordinates are used only for visualization and are not used to calculate the clustering metrics.


\section{Results}
\label{sec:results}

\subsection{Fixed Descriptor Intervals Define Chemically Distinct but Unequal Regions}

Descriptor assignments were recalculated with RDKit 2022.09.5 and matched the stored pretraining assignments for all 250,000 molecules, with zero TPSA mismatches and zero MolLogP mismatches. The 250,000 indexed preprocessing files also cover every source-row index exactly once. Figure~\ref{fig:descriptor-profiles} and Table~\ref{tab:descriptor-regions} show the occupancy and chemical composition of the fixed TPSA and MolLogP regions. The equal-frequency bucketing used during earlier descriptor screening did not generate the assignment arrays used for pretraining.

The TPSA partition is concentrated in the central intervals. Regions 2, 3, and 4, which span 40 to 100 \(\text{\AA}^2\), contain 23.63, 26.74, and 21.32\% of the corpus. Region 0 contains 932 molecules, while region 7 contains 2,708 molecules. Mean hydrogen-bond donor count increases from 0.35 in region 0 to 3.61 in region 7. Mean hydrogen-bond acceptor count increases from 2.52 to 6.98. Molecular weight and heteroatom fraction also increase across most of the TPSA axis. The region profiles are consistent with increasing polar surface and hydrogen-bonding functionality.

MolLogP produces a different occupancy pattern. Regions 0 to 3 and region 5 each contain approximately 14.3\% of the corpus, while region 4 contains 28.37\%. Region 6 contains one molecule and region 7 is empty. Donor and acceptor counts decrease across the populated MolLogP regions, while aromatic atom fraction increases. The two high-lipophilicity branches receive essentially no pretraining examples. Their stored parameters therefore do not represent functions learned from a substantial high-MolLogP population.

\begin{figure*}[t]
\centering
\includegraphics[width=\textwidth]{fig02_descriptor_space_profiles.pdf}
\caption{Fixed TPSA and MolLogP partitions in the ZINC15 250k pretraining corpus. Region occupancy and standardized chemical profiles use the complete corpus. The joint distribution and descriptor correlations use a deterministic 10\% sample for visualization.}
\label{fig:descriptor-profiles}
\end{figure*}

\begin{table*}[t]
\centering
\caption{Occupancy and selected chemical characteristics of the fixed descriptor intervals in the ZINC15 250k pretraining corpus. Values in the final four columns are region means.}
\label{tab:descriptor-regions}
\resizebox{\textwidth}{!}{%
\begin{tabular}{lllrrrrrrr}
\toprule
Descriptor & Region & Interval & Molecules & Fraction (\%) & MW & HBD & HBA & RotB & Rings \\
\midrule
TPSA & R0 & $< 20\ \text{\AA}^{2}$ & 932 & 0.37 & 282.1 & 0.35 & 2.52 & 3.98 & 2.61 \\
 & R1 & $[20, 40)\ \text{\AA}^{2}$ & 17,539 & 7.02 & 294.7 & 0.53 & 2.82 & 4.57 & 2.53 \\
 & R2 & $[40, 60)\ \text{\AA}^{2}$ & 59,080 & 23.63 & 307.0 & 0.97 & 3.47 & 5.02 & 2.48 \\
 & R3 & $[60, 80)\ \text{\AA}^{2}$ & 66,847 & 26.74 & 311.9 & 1.40 & 4.30 & 5.18 & 2.42 \\
 & R4 & $[80, 100)\ \text{\AA}^{2}$ & 53,305 & 21.32 & 313.9 & 1.81 & 5.14 & 5.31 & 2.32 \\
 & R5 & $[100, 120)\ \text{\AA}^{2}$ & 34,637 & 13.85 & 317.0 & 2.31 & 5.86 & 5.46 & 2.25 \\
 & R6 & $[120, 140)\ \text{\AA}^{2}$ & 14,952 & 5.98 & 321.0 & 2.94 & 6.34 & 5.59 & 2.16 \\
 & R7 & $\ge 140\ \text{\AA}^{2}$ & 2,708 & 1.08 & 325.7 & 3.61 & 6.98 & 5.45 & 2.22 \\
\midrule
\MolLogP & R0 & $< -1$ & 35,681 & 14.27 & 311.7 & 2.46 & 5.96 & 5.31 & 2.11 \\
 & R1 & $[-1, 0)$ & 35,759 & 14.30 & 311.3 & 2.00 & 5.47 & 5.29 & 2.26 \\
 & R2 & $[0, 1)$ & 35,986 & 14.39 & 311.1 & 1.67 & 4.97 & 5.29 & 2.35 \\
 & R3 & $[1, 2)$ & 35,968 & 14.39 & 310.9 & 1.40 & 4.41 & 5.20 & 2.44 \\
 & R4 & $[2, 3)$ & 70,922 & 28.37 & 311.1 & 1.19 & 3.79 & 5.10 & 2.49 \\
 & R5 & $[3, 4)$ & 35,683 & 14.27 & 311.3 & 1.02 & 3.35 & 5.01 & 2.52 \\
 & R6 & $[4, 5)$ & 1 & 0.00 & 315.4 & 1.00 & 5.00 & 8.00 & 2.00 \\
 & R7 & $\ge 5$ & 0 & 0.00 &  &  &  &  &  \\
\bottomrule
\end{tabular}}
\end{table*}

\subsection{Historical Placement Screening Favored Final-Layer Conditioning}

The architecture selection stage compares final-layer, odd-layer, and even-layer conditioning with the MolLogP partition and MLP branches. The final-layer model reaches a macro ROC-AUC of 0.7810. Odd-layer and even-layer placement reach 0.7754 and 0.7729, respectively. Final-layer conditioning exceeds each four-layer variant on six of the eight downstream datasets, although no placement is best for every endpoint.

The resulting architecture retains shared feed-forward functions in the first seven graph Transformer layers and introduces descriptor-conditioned computation only in layer 8. This arrangement allows local atom and bond representations to be learned with shared parameters before the descriptor partition modifies the final nonlinear transformation. It also limits the number of conditioned layers and reduces the training-time increase observed for the odd-layer and even-layer variants. All subsequent branch-function experiments use final-layer conditioning. This was an exploratory design screen under the historical downstream protocol. The same test summaries were available during model development, so these values are reported as design provenance rather than an untouched final test comparison.

\begin{figure}[t]
\centering
\includegraphics[width=\columnwidth]{fig03_layer_placement.pdf}
\caption{Placement of descriptor-conditioned computation. The left panel reports five-fold dataset means with sample standard deviations. The right panel reports the eight-dataset macro ROC-AUC and the difference from final-layer placement.}
\label{fig:layer-placement}
\end{figure}

\begin{table}[t]
\centering
\caption{Historical layer-placement screening. Macro ROC-AUC gives equal weight to the eight classification datasets. Dataset-level values are provided in the Supporting Information.}
\label{tab:layer-placement}
\begin{tabular}{lrrrr}
\toprule
Placement & Data sets & Macro AUC & Delta vs last (pp) & Wins vs last \\
\midrule
Last & 8 & 0.7810 & +0.000 & 0/8 \\
Odd & 8 & 0.7754 & -0.564 & 2/8 \\
Even & 8 & 0.7729 & -0.806 & 2/8 \\
\bottomrule
\end{tabular}
\end{table}

\subsection{Historical Branch Screening Produced Descriptor-Specific Function Choices}

After final-layer placement is fixed, each branch is replaced by KAN in a separate pretraining run. Table~\ref{tab:kan-selection} and Figure~\ref{fig:kan-screening} report the downstream change relative to the matched all-MLP model.

For the TPSA partition, regions 0, 1, 2, 3, and 5 have positive eight-dataset macro changes of 0.372, 0.082, 0.578, 0.240, and 0.099 AUC percentage points. Regions 4, 6, and 7 retain MLP. Region 2, which covers 40 to 60 \(\text{\AA}^2\), produces the largest TPSA macro improvement and wins on five datasets. Region 4 also wins on five datasets, but its larger negative changes on other endpoints reduce the macro value below the all-MLP reference.

For the MolLogP partition, regions 1, 2, 3, and 4 have positive macro changes of 0.651, 0.263, 0.810, and 0.163 AUC percentage points. Regions 0, 5, 6, and 7 retain MLP. Region 3, which spans MolLogP values from 1 to 2, gives the largest macro improvement and wins on five datasets. The region-5 screen contains four folds and has a negative macro change, so it remains MLP.

The selected branches are not monotonic along either descriptor axis. Each interval contains diverse scaffolds and correlated molecular properties. The screening results describe branch choices under the historical development protocol rather than a simple chemical rule linking KAN to increasing TPSA or MolLogP. As in the placement study, test summaries were visible during exploratory selection. The selected mixed configurations were subsequently pretrained as complete models from random initialization, while the numerical screen remains design provenance rather than an independent test of branch selection.

\begin{figure*}[t]
\centering
\includegraphics[width=\textwidth]{fig04_kan_selection.pdf}
\caption{Single-branch KAN screening for TPSA and MolLogP. Bars report the eight-dataset macro change relative to the matched all-MLP model. The overlaid curves report pretraining-corpus occupancy, and green bars identify the positive macro changes used for branch selection.}
\label{fig:kan-screening}
\end{figure*}

\begin{table*}[t]
\centering
\caption{Historical single-branch KAN screening after fixing the last-layer placement. During exploratory development, a branch was assigned KAN when its eight-task macro ROC-AUC exceeded the matched all-MLP branch baseline. Delta is reported in percentage points. These values document model construction and are not treated as an independent final test.}
\label{tab:kan-selection}
\resizebox{\textwidth}{!}{%
\begin{tabular}{llrrrrrl}
\toprule
Descriptor & Branch & KAN macro & MLP macro & Delta & Wins & Min. folds & Assignment \\
\midrule
TPSA & k0 & 0.7856 & 0.7818 & +0.372 & 3/8 & 5 & KAN \\
 & k1 & 0.7827 & 0.7818 & +0.082 & 3/8 & 5 & KAN \\
 & k2 & 0.7876 & 0.7818 & +0.578 & 5/8 & 5 & KAN \\
 & k3 & 0.7842 & 0.7818 & +0.240 & 3/8 & 5 & KAN \\
 & k4 & 0.7805 & 0.7818 & -0.131 & 5/8 & 5 & MLP \\
 & k5 & 0.7828 & 0.7818 & +0.099 & 3/8 & 5 & KAN \\
 & k6 & 0.7768 & 0.7818 & -0.500 & 2/8 & 5 & MLP \\
 & k7 & 0.7809 & 0.7818 & -0.092 & 2/8 & 5 & MLP \\
\midrule
\MolLogP & k0 & 0.7776 & 0.7831 & -0.546 & 4/8 & 5 & MLP \\
 & k1 & 0.7896 & 0.7831 & +0.651 & 4/8 & 5 & KAN \\
 & k2 & 0.7857 & 0.7831 & +0.263 & 4/8 & 5 & KAN \\
 & k3 & 0.7912 & 0.7831 & +0.810 & 5/8 & 5 & KAN \\
 & k4 & 0.7847 & 0.7831 & +0.163 & 5/8 & 5 & KAN \\
 & k5 & 0.7829 & 0.7831 & -0.021 & 2/8 & 4 & MLP \\
 & k6 & 0.7792 & 0.7831 & -0.385 & 1/8 & 5 & MLP \\
 & k7 & 0.7789 & 0.7831 & -0.420 & 2/8 & 5 & MLP \\
\bottomrule
\end{tabular}}
\end{table*}

\subsection{Legacy Eight-Layer Snapshot Shows Task-Dependent Predictive Transfer}

Table~\ref{tab:final-performance} reports the legacy eight-layer-aligned development snapshot. All 88 model and dataset combinations contain five scaffold-fold CSV files, giving 440 available files. This file-level completeness does not establish full split coverage because the historical validation loader used random shuffling with \texttt{drop\_last=True} and the historical test loader used \texttt{drop\_last=True}. The checkpoint selected under that validation protocol also cannot be reconstructed because only one best downstream checkpoint was retained. The values below are therefore retained as a development snapshot pending a validation-complete rerun. Within this snapshot, the matched local baseline reaches a macro ROC-AUC of 0.7985.

For TPSA, the all-MLP descriptor-conditioned model reaches a macro ROC-AUC of 0.7924. The selected-KAN model reaches 0.7925. Adding ConLoss to the all-MLP model gives 0.7983. The full TPSA model reaches 0.7810, and the all-KAN control reaches 0.7892. For MolLogP, the all-MLP, selected-KAN, all-MLP with ConLoss, full, and all-KAN models reach 0.7871, 0.7869, 0.7871, 0.7823, and 0.7824.

Dataset-level responses vary across endpoints. TPSA all-MLP conditioning gives 0.8985 on ClinTox, compared with 0.8510 for Base, but is lower on the other seven dataset means. Adding ConLoss to TPSA all-MLP conditioning increases Tox21, ToxCast, SIDER, BACE, HIV, and MUV relative to I1. The MolLogP selected-KAN model is higher than its all-MLP counterpart on Tox21, ToxCast, BACE, HIV, and MUV, while ClinTox decreases from 0.8790 to 0.8356. The endpoint-specific pattern is not captured by the macro average alone.

Published MoleSG results are reported separately in Table~\ref{tab:molesg-comparison}. The published protocol summarizes three runs. The local snapshot uses five scaffold seeds, explicit eight-layer checkpoint loading, and the historical downstream evaluation implementation. The two protocols are displayed separately and are not pooled. Within-protocol descriptive comparisons use the matched local baseline.

\begin{table*}[t]
\centering
\caption{Legacy eight-layer-aligned downstream snapshot. Values are means $\pm$ sample SD across five available scaffold-fold CSV files. Macro values are equal-weight means of the eight dataset means. Historical validation and test loading dropped incomplete batches, so the table is retained as a development result rather than the corrected final benchmark.}
\label{tab:final-performance}
\resizebox{\textwidth}{!}{%
\begin{tabular}{lrrrrrrrrr}
\toprule
Model role & BBBP & Tox21 & ToxCast & SIDER & ClinTox & BACE & HIV & MUV & Macro \\
\midrule
\multicolumn{10}{l}{\textbf{TPSA}} \\
Base & 0.948 $\pm$ 0.018 & 0.816 $\pm$ 0.020 & 0.696 $\pm$ 0.014 & 0.631 $\pm$ 0.026 & 0.851 $\pm$ 0.046 & 0.878 $\pm$ 0.049 & 0.796 $\pm$ 0.032 & 0.772 $\pm$ 0.045 & 0.799 \\
I1 & 0.945 $\pm$ 0.022 & 0.804 $\pm$ 0.020 & 0.692 $\pm$ 0.023 & 0.623 $\pm$ 0.027 & 0.898 $\pm$ 0.044 & 0.871 $\pm$ 0.037 & 0.772 $\pm$ 0.023 & 0.734 $\pm$ 0.053 & 0.792 \\
I1+I2 & 0.944 $\pm$ 0.025 & 0.810 $\pm$ 0.016 & 0.687 $\pm$ 0.016 & 0.616 $\pm$ 0.040 & 0.882 $\pm$ 0.033 & 0.876 $\pm$ 0.046 & 0.785 $\pm$ 0.030 & 0.740 $\pm$ 0.041 & 0.792 \\
I1+I3 & 0.944 $\pm$ 0.026 & 0.811 $\pm$ 0.011 & 0.699 $\pm$ 0.013 & 0.627 $\pm$ 0.037 & 0.870 $\pm$ 0.050 & 0.878 $\pm$ 0.038 & 0.804 $\pm$ 0.041 & 0.755 $\pm$ 0.033 & 0.798 \\
Full & 0.940 $\pm$ 0.019 & 0.810 $\pm$ 0.017 & 0.688 $\pm$ 0.022 & 0.612 $\pm$ 0.034 & 0.862 $\pm$ 0.033 & 0.844 $\pm$ 0.048 & 0.785 $\pm$ 0.036 & 0.707 $\pm$ 0.034 & 0.781 \\
all-KAN & 0.937 $\pm$ 0.034 & 0.798 $\pm$ 0.024 & 0.689 $\pm$ 0.030 & 0.615 $\pm$ 0.029 & 0.887 $\pm$ 0.030 & 0.869 $\pm$ 0.038 & 0.780 $\pm$ 0.029 & 0.739 $\pm$ 0.035 & 0.789 \\
\midrule
\multicolumn{10}{l}{\textbf{\MolLogP}} \\
Base & 0.948 $\pm$ 0.018 & 0.816 $\pm$ 0.020 & 0.696 $\pm$ 0.014 & 0.631 $\pm$ 0.026 & 0.851 $\pm$ 0.046 & 0.878 $\pm$ 0.049 & 0.796 $\pm$ 0.032 & 0.772 $\pm$ 0.045 & 0.799 \\
I1 & 0.930 $\pm$ 0.026 & 0.806 $\pm$ 0.015 & 0.692 $\pm$ 0.024 & 0.623 $\pm$ 0.014 & 0.879 $\pm$ 0.028 & 0.848 $\pm$ 0.034 & 0.782 $\pm$ 0.038 & 0.737 $\pm$ 0.019 & 0.787 \\
I1+I2 & 0.938 $\pm$ 0.008 & 0.815 $\pm$ 0.015 & 0.694 $\pm$ 0.015 & 0.617 $\pm$ 0.029 & 0.836 $\pm$ 0.020 & 0.865 $\pm$ 0.040 & 0.781 $\pm$ 0.050 & 0.750 $\pm$ 0.020 & 0.787 \\
I1+I3 & 0.935 $\pm$ 0.018 & 0.807 $\pm$ 0.013 & 0.685 $\pm$ 0.021 & 0.614 $\pm$ 0.035 & 0.865 $\pm$ 0.041 & 0.855 $\pm$ 0.023 & 0.789 $\pm$ 0.039 & 0.747 $\pm$ 0.021 & 0.787 \\
Full & 0.947 $\pm$ 0.023 & 0.805 $\pm$ 0.017 & 0.691 $\pm$ 0.015 & 0.622 $\pm$ 0.024 & 0.845 $\pm$ 0.019 & 0.849 $\pm$ 0.030 & 0.781 $\pm$ 0.049 & 0.719 $\pm$ 0.041 & 0.782 \\
all-KAN & 0.930 $\pm$ 0.025 & 0.796 $\pm$ 0.018 & 0.693 $\pm$ 0.015 & 0.615 $\pm$ 0.031 & 0.894 $\pm$ 0.017 & 0.851 $\pm$ 0.057 & 0.760 $\pm$ 0.040 & 0.718 $\pm$ 0.041 & 0.782 \\
\bottomrule
\end{tabular}}
\end{table*}

\begin{table}[t]
\centering
\caption{Published MoleSG and local N8 Base results. The published values summarize three runs. The local values summarize five scaffold seeds. They are displayed as separate protocols and are not pooled.}
\label{tab:molesg-comparison}
\begin{tabular}{lrr}
\toprule
Dataset & Published MoleSG & matched local baseline \\
\midrule
BBBP & 0.979 $\pm$ 0.003 & 0.948 $\pm$ 0.018 \\
Tox21 & 0.850 $\pm$ 0.012 & 0.816 $\pm$ 0.020 \\
ToxCast & 0.742 $\pm$ 0.005 & 0.696 $\pm$ 0.014 \\
SIDER & 0.700 $\pm$ 0.002 & 0.631 $\pm$ 0.026 \\
ClinTox & 0.991 $\pm$ 0.009 & 0.851 $\pm$ 0.046 \\
BACE & 0.951 $\pm$ 0.021 & 0.878 $\pm$ 0.049 \\
HIV & 0.877 $\pm$ 0.019 & 0.796 $\pm$ 0.032 \\
MUV & 0.851 $\pm$ 0.008 & 0.772 $\pm$ 0.045 \\
\bottomrule
\end{tabular}
\end{table}

\subsection{Conditional Ablation Shows Nonuniform Component Effects}

Table~\ref{tab:conditional-ablation} and Figure~\ref{fig:conditional-ablation} separate descriptor-conditioned branches from partial KAN and ConLoss within those branches. TPSA I1 is 0.61 AUC percentage points below Base and wins on one of eight datasets. Its 95\% hierarchical bootstrap interval is \(-2.20\) to 1.24 points. Partial KAN changes I1 by 0.01 points and wins on four datasets. ConLoss changes I1 by 0.59 points and wins on six datasets, with an interval from \(-0.86\) to 1.85 points.

The full TPSA model is 1.76 points below Base, with an interval from \(-3.74\) to \(-0.29\) points. The estimated interaction between partial KAN and ConLoss is \(-1.74\) points, with an interval from \(-3.40\) to \(-0.16\) points. The combined configuration is lower than the value implied by adding the two conditional main effects.

For MolLogP, I1 is 1.15 points below Base and wins on one dataset. Its bootstrap interval is \(-2.70\) to 0.46 points. Partial KAN changes I1 by \(-0.01\) points and wins on five datasets. ConLoss changes I1 by 0.00 points and also wins on five datasets. The full MolLogP model is 1.62 points below Base, with an interval from \(-3.42\) to \(-0.26\) points. The estimated interaction is \(-0.46\) points and its interval includes zero.

Selected KAN is also compared with the all-KAN control. The TPSA selected-KAN model is 0.32 points higher than all-KAN and wins on six datasets. The MolLogP selected-KAN model is 0.45 points higher and wins on seven datasets. Both bootstrap intervals include zero. Branch-wise selection is favorable in the observed macro values but is not separated from all-KAN by the reported intervals.

\begin{figure*}[t]
\centering
\includegraphics[width=\textwidth]{fig05_final_n8_performance.pdf}
\caption{Legacy eight-layer-aligned downstream snapshot. Points report means across five available fold files with sample standard deviations, and the lower panels report matched differences from the local Base under the same historical protocol.}
\label{fig:final-performance}
\end{figure*}

\begin{table}[t]
\centering
\caption{Conditional ablation contrasts in the legacy eight-layer-aligned snapshot. Macro deltas and hierarchical bootstrap intervals are reported in ROC-AUC percentage points. The interval resamples datasets and matched folds 10,000 times. These intervals describe the historical files and do not repair incomplete validation or test batches.}
\label{tab:conditional-ablation}
\begin{tabular}{llrrrr}
\toprule
Descriptor & Contrast & Data sets & Delta & Wins & 95\% CI \\
\midrule
TPSA & $\Delta I_1$ & 8 & -0.609 & 1/8 & [-2.204, +1.240] \\
TPSA & $\Delta I_2$ & 8 & +0.007 & 4/8 & [-0.842, +0.832] \\
TPSA & $\Delta I_3$ & 8 & +0.587 & 6/8 & [-0.857, +1.855] \\
TPSA & Full vs Base & 8 & -1.756 & 1/8 & [-3.742, -0.286] \\
MolLogP & $\Delta I_1$ & 8 & -1.146 & 1/8 & [-2.704, +0.455] \\
MolLogP & $\Delta I_2$ & 8 & -0.013 & 5/8 & [-1.540, +1.214] \\
MolLogP & $\Delta I_3$ & 8 & +0.001 & 5/8 & [-0.995, +0.836] \\
MolLogP & Full vs Base & 8 & -1.619 & 0/8 & [-3.417, -0.256] \\
TPSA & $\Delta_{\mathrm{interaction}}$ & 8 & -1.740 &  & [-3.395, -0.158] \\
MolLogP & $\Delta_{\mathrm{interaction}}$ & 8 & -0.461 &  & [-2.045, +1.061] \\
\bottomrule
\end{tabular}
\end{table}


\begin{figure*}[t]
\centering
\includegraphics[width=\textwidth]{fig06_conditional_ablation.pdf}
\caption{Conditional component effects in the legacy eight-layer-aligned snapshot.
Points are matched macro AUC differences and error bars are 95\%
hierarchical bootstrap intervals over datasets and scaffold folds.}
\label{fig:conditional-ablation}
\end{figure*}

\subsection{ConLoss Reorganizes Embeddings along Descriptor Regions}

The representative main-text embedding analysis uses 768 molecules shared across the existing Tox21 fold-1 files. For TPSA, the all-MLP descriptor-conditioned model has a silhouette score of 0.033 and a between-to-within region distance ratio of 1.095. Adding ConLoss increases these values to 0.367 and 1.641. The Davies-Bouldin index decreases from 3.988 to 1.210. The selected-KAN model without ConLoss has a silhouette score of 0.045 and a distance ratio of 1.130, while the full model reaches 0.411 and 1.762.

MolLogP shows the same directional effect. The silhouette score increases from 0.047 for the all-MLP model to 0.336 after ConLoss is added. The distance ratio increases from 1.143 to 1.641. For selected KAN, ConLoss changes the silhouette score from 0.065 to 0.282 and the distance ratio from 1.162 to 1.508. Figure~\ref{fig:embedding-tsne} visualizes the same molecular set. Quantitative metrics are calculated in PCA space rather than from the two-dimensional t-SNE coordinates. The Supporting Information reports the same fixed analysis for all eight datasets and summarizes only datasets with complete five-fold embedding coverage.

The Tox21 set also reveals a mismatch between pretraining and downstream region occupancy. It contains 81 MolLogP region-6 molecules and 101 region-7 molecules, whereas the pretraining corpus contains one and zero molecules in these regions. The associated branches therefore receive little or no direct pretraining support for the high-MolLogP downstream molecules.

ConLoss increases alignment between molecular embeddings and the fixed descriptor regions. The corresponding ROC-AUC effect is smaller and varies by endpoint. Descriptor-region separation and downstream discrimination therefore measure different properties of the learned representation.

\begin{figure*}[t]
\centering
\includegraphics[width=\textwidth]{fig07_embedding_tsne_tox21.pdf}
\caption{t-SNE visualization of Tox21 fold-1 molecular embeddings. Each panel contains the same 768 molecules, and color denotes the descriptor region. Quantitative comparisons use the PCA-space metrics in Table~\ref{tab:embedding-metrics}.}
\label{fig:embedding-tsne}
\end{figure*}

\begin{table}[t]
\centering
\caption{Descriptor-region separation in existing Tox21 fold-1 molecular embeddings. Distances and clustering indices are calculated after model-specific PCA and L2 normalization.}
\label{tab:embedding-metrics}
\begin{tabular}{llrrrrr}
\toprule
Partition & Model & Silhouette & DB index & Intra & Inter & Ratio \\
\midrule
TPSA & I1 & 0.033 & 3.988 & 1.296 & 1.419 & 1.095 \\
TPSA & I1+I2 & 0.045 & 3.181 & 1.262 & 1.426 & 1.130 \\
TPSA & I1+I3 & 0.367 & 1.210 & 0.906 & 1.487 & 1.641 \\
TPSA & Full & 0.411 & 1.081 & 0.848 & 1.494 & 1.762 \\
\MolLogP & I1 & 0.047 & 3.611 & 1.248 & 1.427 & 1.143 \\
\MolLogP & I1+I2 & 0.065 & 3.244 & 1.231 & 1.430 & 1.162 \\
\MolLogP & I1+I3 & 0.336 & 1.279 & 0.900 & 1.477 & 1.641 \\
\MolLogP & Full & 0.282 & 1.484 & 0.974 & 1.469 & 1.508 \\
\bottomrule
\end{tabular}
\end{table}

\subsection{Descriptor Conditioning Increases Stored Capacity and Runtime}

Table~\ref{tab:efficiency} separates stored parameters from active feed-forward arithmetic. The eight dense feed-forward layers contain 1.053 million parameters and require an analytical 2.097 million FLOPs per token. Replacing the final dense feed-forward network with eight stored MLP branches increases stored feed-forward parameters to 1.974 million. One active MLP branch retains the same analytical feed-forward FLOP count, and descriptor assignment adds no learned router parameters.

KAN branches increase both storage and active arithmetic. The TPSA and MolLogP selected-KAN models store 7.214 and 6.166 million feed-forward parameters. Active FLOP estimates weighted by region occupancy exceed the dense model because each KAN branch evaluates a base mapping and spline coefficients. The all-KAN control stores 10.358 million feed-forward parameters and has a lower-bound active count of 4.194 million FLOPs per token. The learned routed top-1 control stores 8.421 million feed-forward parameters and 16,384 router parameters. Its active count is 2.130 million FLOPs per token before dispatch and softmax overhead.

Observed epoch time follows the same pattern. After segment startup and timing outliers are excluded, the matched local baseline requires a median of 908.8 s per active pretraining epoch. TPSA and MolLogP all-MLP final-layer models require 918.7 and 951.5 s. Their selected-KAN models require 1032.1 and 1032.6 s. The full TPSA and MolLogP models require 1032.1 and 1007.0 s. Historical odd-layer and even-layer conditioning each require 1035.6 s per epoch, compared with 908.6 s for final-layer conditioning. In the present implementation, fixed assignment removes the learned router but does not reduce wall-clock time below the dense baseline.

Slurm accounting for downstream jobs is reported in the Supporting Information. It records completed and unsuccessful allocation time separately. These values describe execution history and are not used as a controlled runtime comparison because retry jobs can skip previously completed folds.



\begin{figure}[H]
\centering
\includegraphics[width=\columnwidth]{fig08_performance_compute.pdf}
\caption{Matched downstream macro ROC-AUC as a function of analytical active feed-forward computation and recovered active pretraining time. Point area denotes analytical stored model parameters. KAN FLOPs are architecture-based lower bounds.}
\label{fig:performance-compute}
\end{figure}

\begin{table*}[t]
\centering
\caption{Analytical model complexity and observed successful-job runtime. Stored totals use the shared non-FFN parameter count plus the implemented branch parameters. Active FFN parameters for mixed models are occupancy-weighted. KAN FLOPs exclude spline-basis evaluation overhead. Finetuning times describe Slurm accounting records and are not a controlled architecture benchmark.}
\label{tab:efficiency}
\resizebox{\textwidth}{!}{%
\begin{tabular}{lrrrrrrr}
\toprule
Configuration & Stored params (M) & Active FFN params (M) & FLOPs/token (M) & PT min/epoch & PT GPU h & FT h/job & FT GPU h \\
\midrule
Vanilla & 3.744 & 1.053 & 2.097 & 15.15 & 76.0 & 4.47 & 53.6 \\
TPSA all-MLP & 4.665 & 1.053 & 2.097 & 15.31 & 77.0 & 1.15 & 71.5 \\
MolLogP all-MLP & 4.665 & 1.053 & 2.097 & 15.86 & 79.5 & 1.18 & 69.1 \\
TPSA selected-KAN & 9.905 & 1.803 & 3.599 & 17.20 & 86.2 & 1.21 & 80.4 \\
MolLogP selected-KAN & 8.857 & 1.802 & 3.596 & 17.21 & 86.1 & 1.30 & 60.6 \\
TPSA Full & 9.905 & 1.803 & 3.599 & 17.20 & 86.2 & 1.16 & 71.2 \\
MolLogP Full & 8.857 & 1.802 & 3.596 & 16.78 & 84.8 & 1.37 & 38.0 \\
TPSA all-KAN & 13.049 & 2.101 & 4.194 & 17.53 & 89.0 & 3.66 & 61.2 \\
MolLogP all-KAN & 13.049 & 2.101 & 4.194 & 17.73 & 91.1 & 1.44 & 76.7 \\
\bottomrule
\end{tabular}}
\end{table*}


\section{Discussion}
\label{sec:discussion}

\subsection{Physicochemical Descriptors as Branch-Assignment Variables}

The central methodological distinction is the role assigned to the descriptor. MolLogP and TPSA are calculated from the same standardized molecular structure provided to the graph encoder. They do not add an experimental measurement to the pretraining data. Their role is to select the final-layer function applied after shared structural encoding. This differs from descriptor concatenation, where all molecules pass through the same network and the descriptor enters the predictor as another input variable.

Fixed assignment makes the molecular population of each branch directly measurable. The TPSA regions show progressive changes in hydrogen-bond donor count, acceptor count, heteroatom fraction, and molecular weight. The populated MolLogP regions show decreasing donor and acceptor counts and increasing aromatic atom fraction. These profiles are consistent with the intended physicochemical coordinates and provide chemical context for the otherwise numerical branch indices.

The same analysis also reveals a limitation. Fixed intervals inherit the empirical distribution of the pretraining corpus. TPSA produces several well-populated central regions and sparse extremes. MolLogP produces one region with a single molecule and one empty region. Those branches add stored parameters without receiving representative pretraining updates. The presence of many high-MolLogP molecules in the downstream Tox21 embedding set further shows that pretraining support and application-domain support can differ. Descriptor-conditioned models therefore require region-occupancy analysis in both domains.

\subsection{Branch Function Selection Is Partition and Protocol Dependent}

TPSA and MolLogP produce different MLP and KAN choices. This result follows from the different molecular populations created by the two descriptor axes. Each interval contains multiple scaffolds and correlated physicochemical properties. A selected KAN branch should therefore be interpreted as a function choice for one region under the current training and evaluation protocol, not as a general statement that a named chemical class requires KAN.

Several screening effects are small. Positive macro differences below 0.4 ROC-AUC percentage points determine four of the five TPSA KAN choices. These choices can be sensitive to scaffold composition and dataset weighting. In the legacy eight-layer snapshot, selected KAN does not consistently improve the matched all-MLP model. The comparison with all-KAN is more favourable, especially for MolLogP, but the reported bootstrap intervals include zero. Branch-wise screening avoided uniform KAN replacement in the implemented models. The present files do not establish a general advantage of KAN over MLP for physicochemical regions.

The separation between architecture selection and final evaluation is therefore important. Layer placement and single-branch KAN screening define the model configuration. The complete selected-KAN models are then pretrained jointly from random initialization. Their performance cannot be inferred by adding the individual screening improvements. A submission-grade evaluation also requires selection on complete validation splits followed by one evaluation on complete test splits. The historical screen did not satisfy that separation and is retained as model-development provenance.

\subsection{Embedding Alignment and Property Prediction Are Different Objectives}

ConLoss produces the clearest representation-level effect. Silhouette scores and between-to-within region distance ratios increase for both descriptor axes, while Davies-Bouldin indices decrease. The molecular embeddings therefore retain the imposed MolLogP or TPSA partition more strongly after contrastive regularization.

The downstream effect is smaller and varies across endpoints. TPSA all-MLP with ConLoss gives the strongest positive conditional macro contrast, but its bootstrap interval includes zero. MolLogP shows almost no macro change after ConLoss is added. Toxicity and bioactivity depend on scaffold, stereochemistry, ionization, local interactions, and assay conditions in addition to global lipophilicity or polar surface. Stronger separation along a descriptor axis can therefore occur without producing a better decision boundary for every endpoint.

The embedding and ROC-AUC analyses answer different questions. Descriptor-region metrics test whether the representation follows the imposed physicochemical coordinate. Downstream evaluation tests whether that organization transfers to the endpoint labels. The results show strong control of representation geometry and limited, task-dependent predictive transfer.

\subsection{Fixed Assignment Does Not Reduce Runtime}

Fixed descriptor assignment removes router logits and load-balancing terms. It does not reduce stored capacity because every branch remains in the model. Replacing one dense feed-forward network with eight MLP branches increases stored parameters, although one active MLP branch retains the same analytical feed-forward FLOP count. Selected KAN increases both stored capacity and active arithmetic.

Observed epoch times follow the same pattern. All-MLP descriptor conditioning is slightly slower than the dense baseline, and selected-KAN configurations are slower than all-MLP conditioning. Mask-based dispatch and spline evaluation remain part of the executed computation. The practical benefit of fixed assignment is therefore the reproducible relation between descriptor value and branch choice rather than lower runtime.

\subsection{Scope and Limitations}

MolLogP and TPSA represent two global properties of standardized two-dimensional structures. MolLogP does not resolve the pH-dependent distribution of ionized species. TPSA does not describe conformer-dependent polar surface exposure or substituent electronic effects. Both depend on molecular standardization, and hard interval boundaries can send neighbouring values to different branches.

The experimental scope includes one pretraining corpus, one multimodal backbone, eight classification datasets, and five scaffold seeds. Placement, branch screening, and the eight-layer snapshot use the historical downstream evaluation code. The model weights are explicitly aligned to an eight-layer backbone, while validation and test coverage remain limited by incomplete final batches. The statistical analysis describes the available files and scaffold partitions rather than a corrected final benchmark, unseen assay families, or prospective measurements.

Other molecular descriptors could be studied within the same computation pattern when their definitions, chemical relevance, and region support are established. That extension requires separate experiments. The present conclusions are limited to fixed MolLogP and TPSA assignments in the reported small-molecule pretraining and downstream tasks.

\section{Conclusions}

This study examines physicochemical descriptors as direct branch-assignment variables during molecular pretraining. Fixed MolLogP and TPSA intervals select one of eight branches in the final layer of an eight-layer graph Transformer. The analysis separates descriptor-region chemistry, layer placement, branch-wise MLP and KAN selection, descriptor-aware contrastive learning, downstream prediction, representation geometry, and computational cost.

Historical placement screening favored final-layer conditioning, and TPSA and MolLogP produced different branch-function choices during single-branch screening. In the legacy eight-layer snapshot, the selected mixed models do not consistently improve the matched all-MLP variants. ConLoss increases descriptor-region separation across both descriptor axes, while its ROC-AUC effect is smaller and task-dependent. The complete TPSA and MolLogP configurations remain below the matched local Base in the same snapshot.

The observations separate representation organization from predictive transfer. A fixed physicochemical coordinate determines the branch used to process a molecular representation and, with ConLoss, produces marked alignment between embeddings and descriptor regions. The same organization does not improve every molecular property endpoint in the historical downstream files. Descriptor-conditioned computation therefore provides a direct setting for studying how global molecular variables interact with learned structural representations, while predictive performance, representation geometry, stored capacity, and runtime remain distinct outcomes. Corrected validation-complete and test-complete downstream evaluation is required before these predictive values are used as submission results.

\section*{Supporting Information}

The Supporting Information contains the complete mathematical formulation,
descriptor-region profiles, dataset-level placement and KAN screens,
fold-level downstream tables, conditional contrasts, auxiliary controls,
embedding analyses for all eight datasets, runtime records, and metric
definitions.

\section*{Data and Code Availability}

The prepared reproducibility package contains training and analysis code,
configuration guidance for placement and branch-function experiments,
fold-level ROC-AUC values, descriptor-region profiles, branch-selection
summaries, embedding metrics, figure source data, and runtime records. The
ZINC15 and MoleculeNet source collections remain subject to their original
distribution terms and are referenced by path rather than redistributed. A
public repository URL and archival identifier will be inserted after the
authors approve release of the reviewer package.

\section*{Acknowledgments}

OpenAI language tools assisted manuscript language editing. GPT Image2 was
used to produce nonquantitative background elements for Figure~\ref{fig:workflow}.
The molecular structure, descriptor values, interval assignments, branch
functions, labels, arrows, and all quantitative figures were generated from
RDKit or project code and checked by the authors. The table-of-contents graphic
contains no generative-image elements.

\ifdevelopmentdraft
\noindent\textit{Author input required before submission includes funding and
grant information, ORCID identifiers, the CRediT contribution statement, the
competing-interest statement, and public repository identifiers.}
\fi

\bibliography{references}

\clearpage
\appendix
\setcounter{section}{0}
\setcounter{table}{0}
\setcounter{figure}{0}
\setcounter{equation}{0}
\renewcommand{\thesection}{S\arabic{section}}
\renewcommand{\thetable}{S\arabic{table}}
\renewcommand{\thefigure}{S\arabic{figure}}
\renewcommand{\theequation}{S\arabic{equation}}
\providecommand{\theHsection}{}
\providecommand{\theHtable}{}
\providecommand{\theHfigure}{}
\providecommand{\theHequation}{}
\renewcommand{\theHsection}{S.\arabic{section}}
\renewcommand{\theHtable}{S.\arabic{table}}
\renewcommand{\theHfigure}{S.\arabic{figure}}
\renewcommand{\theHequation}{S.\arabic{equation}}

\section*{Integrated Supporting Information}

This appendix reports the numerical and visualization record underlying the
main-text development snapshot. Historical design-selection experiments are
kept separate from the legacy eight-layer-aligned evaluation. That evaluation
contains all 440 planned fold CSV files, and the N4 depth reference contains
all 40 planned fold CSV files. File completeness does not imply complete
validation or test sample coverage under the historical \texttt{drop\_last}
settings.

\section{Mathematical Formulation of Descriptor-Guided Computation}
\label{sec:math-appendix}

\subsection{Descriptor-Defined Regions}

Let \(\mathcal{M}\) denote the molecular chemical space represented in the
pretraining corpus. A scalar molecular descriptor is a deterministic map
\begin{equation}
d:\mathcal{M}\rightarrow\mathbb{R}.
\end{equation}
For molecule \(m_i\), its coordinate on the selected descriptor axis is
\begin{equation}
z_i=d(m_i).
\end{equation}
The descriptor axis is divided into \(K\) fixed intervals
\begin{equation}
\mathcal{P}_{d}=\{I_1,I_2,\ldots,I_K\}.
\end{equation}
The region assignment is
\begin{equation}
e_i=\phi(z_i), \qquad
\phi(z_i)=k \quad \text{if } z_i\in I_k.
\label{eq:descriptor-assignment}
\end{equation}
The mapping \(\phi\) is fixed before model training. It contains no learned
router weights and does not use downstream labels.

For the experiments in this work, \(K=8\). The TPSA boundaries are 20, 40,
60, 80, 100, 120, and \(140\ \text{\AA}^{2}\). The MolLogP boundaries are
\(-1\), 0, 1, 2, 3, 4, and 5. These are fixed chemical intervals rather than
empirical quantiles.

\subsection{Descriptor-Conditioned Computation Branch}

Let \(\mathbf{H}^{(\ell)}_i\in\mathbb{R}^{n_i\times d_h}\) denote the node
representations of molecule \(m_i\) at Transformer layer \(\ell\), where
\(n_i\) is its number of graph nodes and \(d_h=256\). At a
descriptor-conditioned layer, every valid node of the molecule is processed
by the branch associated with \(e_i\).
\begin{equation}
\widetilde{\mathbf{H}}^{(\ell)}_i
=E^{(\ell)}_{e_i}\left(\mathbf{H}^{(\ell)}_i\right).
\label{eq:conditional-branch}
\end{equation}
The branch output then enters the residual and normalization operations of
the original graph Transformer layer. One branch is active for each molecule.
All branches are stored in the model.

For an all-MLP configuration, branch \(k\) applies
\begin{equation}
E^{\mathrm{MLP}}_k(\mathbf{x})
=
\mathbf{W}_{k,2}
\operatorname{Mish}
\left(\mathbf{W}_{k,1}\mathbf{x}+\mathbf{b}_{k,1}\right)
+\mathbf{b}_{k,2}.
\end{equation}
Both linear transformations map between 256-dimensional spaces in the
current implementation.

\subsection{Layer Placement}

The graph encoder contains \(L=8\) Transformer layers. The placement study
uses the following one-indexed layer sets.
\begin{align}
\mathcal{L}_{\mathrm{last}} &= \{8\},\\
\mathcal{L}_{\mathrm{odd}} &= \{1,3,5,7\},\\
\mathcal{L}_{\mathrm{even}} &= \{2,4,6,8\}.
\end{align}
The final architecture uses \(\mathcal{L}_{\mathrm{last}}\). Layers 1 through
7 therefore retain shared dense FFNs, and descriptor-conditioned computation
is introduced only at layer 8.

\subsection{Partial KAN Function Selection}

The KAN branch is composed of two KAN linear transformations. Each
transformation combines a SiLU base mapping with a cubic B-spline mapping.
\begin{equation}
\mathcal{K}(\mathbf{x})
=
\mathbf{W}_{b}\operatorname{SiLU}(\mathbf{x})
+
\mathbf{W}_{s}\mathbf{B}_{3}(\mathbf{x}),
\end{equation}
where \(\mathbf{B}_{3}\) denotes the order-3 B-spline basis evaluated on a
grid of size 5. The branch function is selected separately for each
descriptor-defined region.
\begin{equation}
E_k=
\begin{cases}
E_k^{\mathrm{KAN}}, & a_k=1,\\
E_k^{\mathrm{MLP}}, & a_k=0.
\end{cases}
\label{eq:partial-kan}
\end{equation}

For branch \(k\), the selection indicator is determined from the matched
single-branch screening experiment.
\begin{equation}
a_k
=
\mathbb{I}
\left[
\frac{1}{|\mathcal{D}|}
\sum_{q\in\mathcal{D}}
\left(
\overline{\operatorname{AUC}}_{q,k}^{\mathrm{KAN}}
-
\overline{\operatorname{AUC}}_{q}^{\mathrm{MLP}}
\right)
>0
\right],
\label{eq:kan-selection}
\end{equation}
where \(\mathcal{D}\) contains the eight downstream classification datasets
and each mean is calculated over the available matched scaffold seeds. All
screening configurations contained five folds except MolLogP branch 5, which
contained a minimum of four folds and was not selected. The locked TPSA KAN
branches are \(0,1,2,3,\) and \(5\). The locked MolLogP KAN branches are
\(1,2,3,\) and \(4\).

\subsection{Descriptor-Aware Contrastive Regularization}

The graph encoder output is mean-pooled over valid nodes.
\begin{equation}
\mathbf{g}_i
=
\frac{\sum_{j=1}^{n_{\max}}m_{ij}\mathbf{h}_{ij}}
{\max\left(1,\sum_{j=1}^{n_{\max}}m_{ij}\right)},
\end{equation}
where \(m_{ij}\) is the node mask. The normalized molecular representation is
\begin{equation}
\mathbf{u}_i=\frac{\mathbf{g}_i}{\lVert\mathbf{g}_i\rVert_2}.
\end{equation}
For anchor \(i\), the positive set contains other molecules in the same
descriptor-defined region.
\begin{equation}
\mathcal{P}(i)=\{p\neq i:e_p=e_i\}.
\end{equation}
The contrastive term implemented in the pretraining code is
\begin{equation}
\mathcal{L}_{\mathrm{Con}}
=
-\frac{1}{|\mathcal{A}|}
\sum_{i\in\mathcal{A}}
\frac{1}{|\mathcal{P}(i)|}
\sum_{p\in\mathcal{P}(i)}
\log
\frac{
\exp\left(\mathbf{u}_i^{\mathsf{T}}\mathbf{u}_p/\tau\right)
}{
\sum_{a\neq i}
\exp\left(\mathbf{u}_i^{\mathsf{T}}\mathbf{u}_a/\tau\right)
},
\label{eq:conloss}
\end{equation}
where \(\mathcal{A}\) contains anchors with at least one positive example and
\(\tau=0.1\). The total pretraining objective is
\begin{equation}
\mathcal{L}
=
\mathcal{L}_{\mathrm{SCE}}
+
\mathcal{L}_{\mathrm{MLM}}
+
\lambda_{\mathrm{Con}}\mathcal{L}_{\mathrm{Con}},
\end{equation}
with \(\lambda_{\mathrm{Con}}=0.1\). SCE denotes scaled cosine error for masked
graph reconstruction. MLM denotes masked language modeling over the tokenized
simplified molecular input line entry system representation.

\subsection{Conditional Ablation Contrasts}

Partial KAN and ConLoss are defined only after descriptor-conditioned
branches are present. Their effects are therefore reported conditionally.
\begin{align}
\Delta I_1 &= I_1-\mathrm{Base},\\
\Delta I_2 &= (I_1+I_2)-I_1,\\
\Delta I_3 &= (I_1+I_3)-I_1.
\end{align}
The interaction between partial KAN and ConLoss is
\begin{equation}
\Delta_{\mathrm{interaction}}
=
\mathrm{Full}
-(I_1+I_2)
-(I_1+I_3)
+I_1.
\end{equation}
These quantities are calculated from matched five-fold dataset means. They
are not interpreted as a complete independent \(2^3\) factorial design.

\section{Descriptor Distributions and Region Chemistry}

The descriptor intervals are fixed before pretraining and are not empirical
quantiles. Figure~\ref{fig:si-descriptor-joint} reports the underlying
one-dimensional distributions and their joint projection. The complete
region-level chemical profiles are provided in
Table~\ref{tab:si-descriptor-profiles}.

\begin{figure*}[t]
\centering
\includegraphics[width=\textwidth]{figS01_descriptor_joint_distribution.pdf}
\caption{TPSA and MolLogP distributions in a deterministic 10\% sample of the
ZINC15 250k pretraining corpus. Vertical lines mark the fixed interval
boundaries. The joint panel shows that the two descriptor axes are related but
not interchangeable.}
\label{fig:si-descriptor-joint}
\end{figure*}

\begin{table*}[t]
\centering
\caption{Complete chemical profiles of the fixed descriptor-defined regions in the ZINC15 250k pretraining corpus. Values after occupancy are region means.}
\label{tab:si-descriptor-profiles}
\resizebox{\textwidth}{!}{%
\begin{tabular}{lllrrrrrrrrrr}
\toprule
Descriptor & Region & Interval & Molecules & Fraction (\%) & MW & HBD & HBA & RotB & Rings & Aromatic frac. & Heteroatom frac. & Charge \\
\midrule
TPSA & R0 & $< 20\ \text{\AA}^{2}$ & 932 & 0.37 & 282.1 & 0.35 & 2.52 & 3.98 & 2.61 & 0.268 & 0.188 & 0.002 \\
 & R1 & $[20, 40)\ \text{\AA}^{2}$ & 17,539 & 7.02 & 294.7 & 0.53 & 2.82 & 4.57 & 2.53 & 0.257 & 0.209 & 0.001 \\
 & R2 & $[40, 60)\ \text{\AA}^{2}$ & 59,080 & 23.63 & 307.0 & 0.97 & 3.47 & 5.02 & 2.48 & 0.273 & 0.240 & 0.000 \\
 & R3 & $[60, 80)\ \text{\AA}^{2}$ & 66,847 & 26.74 & 311.9 & 1.40 & 4.30 & 5.18 & 2.42 & 0.281 & 0.283 & 0.000 \\
 & R4 & $[80, 100)\ \text{\AA}^{2}$ & 53,305 & 21.32 & 313.9 & 1.81 & 5.14 & 5.31 & 2.32 & 0.259 & 0.335 & 0.000 \\
 & R5 & $[100, 120)\ \text{\AA}^{2}$ & 34,637 & 13.85 & 317.0 & 2.31 & 5.86 & 5.46 & 2.25 & 0.254 & 0.381 & 0.000 \\
 & R6 & $[120, 140)\ \text{\AA}^{2}$ & 14,952 & 5.98 & 321.0 & 2.94 & 6.34 & 5.59 & 2.16 & 0.267 & 0.414 & 0.000 \\
 & R7 & $\ge 140\ \text{\AA}^{2}$ & 2,708 & 1.08 & 325.7 & 3.61 & 6.98 & 5.45 & 2.22 & 0.304 & 0.446 & 0.000 \\
\midrule
\MolLogP & R0 & $< -1$ & 35,681 & 14.27 & 311.7 & 2.46 & 5.96 & 5.31 & 2.11 & 0.166 & 0.403 & 0.000 \\
 & R1 & $[-1, 0)$ & 35,759 & 14.30 & 311.3 & 2.00 & 5.47 & 5.29 & 2.26 & 0.205 & 0.362 & 0.000 \\
 & R2 & $[0, 1)$ & 35,986 & 14.39 & 311.1 & 1.67 & 4.97 & 5.29 & 2.35 & 0.240 & 0.326 & 0.000 \\
 & R3 & $[1, 2)$ & 35,968 & 14.39 & 310.9 & 1.40 & 4.41 & 5.20 & 2.44 & 0.278 & 0.289 & 0.000 \\
 & R4 & $[2, 3)$ & 70,922 & 28.37 & 311.1 & 1.19 & 3.79 & 5.10 & 2.49 & 0.315 & 0.252 & 0.000 \\
 & R5 & $[3, 4)$ & 35,683 & 14.27 & 311.3 & 1.02 & 3.35 & 5.01 & 2.52 & 0.362 & 0.226 & 0.000 \\
 & R6 & $[4, 5)$ & 1 & 0.00 & 315.4 & 1.00 & 5.00 & 8.00 & 2.00 & 0.522 & 0.217 & 0.000 \\
 & R7 & $\ge 5$ & 0 & 0.00 &  &  &  &  &  &  &  &  \\
\bottomrule
\end{tabular}}
\end{table*}

\section{Layer-Placement Screening}

The placement experiment asks where descriptor-conditioned computation should
be introduced. Figure~\ref{fig:si-layer-heatmap} and
Table~\ref{tab:si-layer-detail} provide the complete dataset-level screening
results. These historical design-selection results are not pooled with the
legacy eight-layer-aligned snapshot.

\begin{figure*}[t]
\centering
\includegraphics[width=\textwidth]{figS02_layer_placement_heatmap.pdf}
\caption{Layer-placement screening across the eight classification datasets.
The left panel reports absolute five-fold ROC-AUC. The right panel reports the
difference from final-layer conditioning in AUC percentage points.}
\label{fig:si-layer-heatmap}
\end{figure*}

\begin{table*}[t]
\centering
\caption{Complete layer-placement screening results. Values are five-fold ROC-AUC mean $\pm$ sample SD under the historical design-selection protocol.}
\label{tab:si-layer-detail}
\resizebox{\textwidth}{!}{%
\begin{tabular}{lrrrrrrrrr}
\toprule
Placement & BBBP & Tox21 & ToxCast & SIDER & ClinTox & BACE & HIV & MUV & Macro \\
\midrule
Last & 0.937 $\pm$ 0.031 & 0.809 $\pm$ 0.017 & 0.690 $\pm$ 0.015 & 0.602 $\pm$ 0.020 & 0.878 $\pm$ 0.039 & 0.831 $\pm$ 0.030 & 0.765 $\pm$ 0.038 & 0.737 $\pm$ 0.045 & 0.7810 \\
Odd & 0.928 $\pm$ 0.015 & 0.783 $\pm$ 0.015 & 0.684 $\pm$ 0.017 & 0.612 $\pm$ 0.019 & 0.886 $\pm$ 0.032 & 0.822 $\pm$ 0.047 & 0.762 $\pm$ 0.036 & 0.726 $\pm$ 0.042 & 0.7754 \\
Even & 0.927 $\pm$ 0.026 & 0.789 $\pm$ 0.014 & 0.677 $\pm$ 0.019 & 0.613 $\pm$ 0.027 & 0.853 $\pm$ 0.077 & 0.820 $\pm$ 0.059 & 0.760 $\pm$ 0.041 & 0.745 $\pm$ 0.041 & 0.7729 \\
\bottomrule
\end{tabular}}
\end{table*}

\section{Single-Branch KAN Screening}

TPSA and MolLogP were screened independently after final-layer conditioning
had been fixed. Figures~\ref{fig:si-tpsa-kan-absolute} to
\ref{fig:si-logp-kan-delta} separate absolute performance from the matched
all-MLP difference. Tables~\ref{tab:si-tpsa-kan-absolute} to
\ref{tab:si-logp-kan-deltas} provide the corresponding numerical values.

\begin{figure*}[t]
\centering
\includegraphics[width=\textwidth]{figS03_tpsa_kan_absolute.pdf}
\caption{Absolute ROC-AUC for the TPSA single-branch KAN screening models.
Rows identify the branch replaced by KAN.}
\label{fig:si-tpsa-kan-absolute}
\end{figure*}

\begin{figure*}[t]
\centering
\includegraphics[width=\textwidth]{figS04_tpsa_kan_delta.pdf}
\caption{TPSA single-branch KAN differences from the matched all-MLP
final-layer model. Values are AUC percentage points.}
\label{fig:si-tpsa-kan-delta}
\end{figure*}

\begin{figure*}[t]
\centering
\includegraphics[width=\textwidth]{figS05_logp_kan_absolute.pdf}
\caption{Absolute ROC-AUC for the MolLogP single-branch KAN screening models.
Rows identify the branch replaced by KAN.}
\label{fig:si-logp-kan-absolute}
\end{figure*}

\begin{figure*}[t]
\centering
\includegraphics[width=\textwidth]{figS06_logp_kan_delta.pdf}
\caption{MolLogP single-branch KAN differences from the matched all-MLP
final-layer model. Values are AUC percentage points.}
\label{fig:si-logp-kan-delta}
\end{figure*}

\begin{table*}[t]
\centering
\caption{TPSA single-branch KAN screening. Values are ROC-AUC mean $\pm$ sample SD. Macro averages give equal weight to the eight dataset means.}
\label{tab:si-tpsa-kan-absolute}
\resizebox{\textwidth}{!}{%
\begin{tabular}{lrrrrrrrrrr}
\toprule
Branch & BBBP & Tox21 & ToxCast & SIDER & ClinTox & BACE & HIV & MUV & Macro & Min. folds \\
\midrule
k0 & 0.935 $\pm$ 0.034 & 0.806 $\pm$ 0.019 & 0.690 $\pm$ 0.021 & 0.616 $\pm$ 0.041 & 0.879 $\pm$ 0.057 & 0.863 $\pm$ 0.051 & 0.787 $\pm$ 0.023 & 0.708 $\pm$ 0.040 & 0.7856 & 5 \\
k1 & 0.927 $\pm$ 0.043 & 0.805 $\pm$ 0.027 & 0.694 $\pm$ 0.020 & 0.614 $\pm$ 0.028 & 0.858 $\pm$ 0.079 & 0.860 $\pm$ 0.072 & 0.780 $\pm$ 0.034 & 0.724 $\pm$ 0.025 & 0.7827 & 5 \\
k2 & 0.925 $\pm$ 0.048 & 0.808 $\pm$ 0.018 & 0.695 $\pm$ 0.016 & 0.618 $\pm$ 0.044 & 0.889 $\pm$ 0.021 & 0.842 $\pm$ 0.049 & 0.785 $\pm$ 0.043 & 0.739 $\pm$ 0.044 & 0.7876 & 5 \\
k3 & 0.927 $\pm$ 0.044 & 0.802 $\pm$ 0.028 & 0.693 $\pm$ 0.015 & 0.616 $\pm$ 0.027 & 0.877 $\pm$ 0.034 & 0.837 $\pm$ 0.057 & 0.780 $\pm$ 0.037 & 0.743 $\pm$ 0.040 & 0.7842 & 5 \\
k4 & 0.941 $\pm$ 0.029 & 0.807 $\pm$ 0.024 & 0.693 $\pm$ 0.019 & 0.614 $\pm$ 0.024 & 0.861 $\pm$ 0.040 & 0.836 $\pm$ 0.057 & 0.762 $\pm$ 0.030 & 0.730 $\pm$ 0.036 & 0.7805 & 5 \\
k5 & 0.939 $\pm$ 0.034 & 0.805 $\pm$ 0.028 & 0.693 $\pm$ 0.017 & 0.616 $\pm$ 0.019 & 0.892 $\pm$ 0.022 & 0.851 $\pm$ 0.068 & 0.775 $\pm$ 0.035 & 0.692 $\pm$ 0.023 & 0.7828 & 5 \\
k6 & 0.927 $\pm$ 0.036 & 0.805 $\pm$ 0.019 & 0.697 $\pm$ 0.018 & 0.613 $\pm$ 0.041 & 0.837 $\pm$ 0.063 & 0.838 $\pm$ 0.061 & 0.775 $\pm$ 0.029 & 0.723 $\pm$ 0.040 & 0.7768 & 5 \\
k7 & 0.931 $\pm$ 0.033 & 0.799 $\pm$ 0.016 & 0.694 $\pm$ 0.014 & 0.618 $\pm$ 0.038 & 0.878 $\pm$ 0.052 & 0.841 $\pm$ 0.050 & 0.770 $\pm$ 0.035 & 0.717 $\pm$ 0.037 & 0.7809 & 5 \\
\bottomrule
\end{tabular}}
\end{table*}
\begin{table*}[t]
\centering
\caption{\MolLogP single-branch KAN screening. Values are ROC-AUC mean $\pm$ sample SD. Macro averages give equal weight to the eight dataset means.}
\label{tab:si-logp-kan-absolute}
\resizebox{\textwidth}{!}{%
\begin{tabular}{lrrrrrrrrrr}
\toprule
Branch & BBBP & Tox21 & ToxCast & SIDER & ClinTox & BACE & HIV & MUV & Macro & Min. folds \\
\midrule
k0 & 0.936 $\pm$ 0.025 & 0.801 $\pm$ 0.026 & 0.690 $\pm$ 0.022 & 0.613 $\pm$ 0.029 & 0.872 $\pm$ 0.054 & 0.822 $\pm$ 0.052 & 0.775 $\pm$ 0.047 & 0.713 $\pm$ 0.025 & 0.7776 & 5 \\
k1 & 0.929 $\pm$ 0.021 & 0.810 $\pm$ 0.010 & 0.697 $\pm$ 0.021 & 0.628 $\pm$ 0.026 & 0.903 $\pm$ 0.037 & 0.832 $\pm$ 0.037 & 0.782 $\pm$ 0.048 & 0.737 $\pm$ 0.032 & 0.7896 & 5 \\
k2 & 0.934 $\pm$ 0.021 & 0.808 $\pm$ 0.011 & 0.689 $\pm$ 0.017 & 0.615 $\pm$ 0.028 & 0.898 $\pm$ 0.047 & 0.840 $\pm$ 0.043 & 0.758 $\pm$ 0.051 & 0.744 $\pm$ 0.023 & 0.7857 & 5 \\
k3 & 0.930 $\pm$ 0.028 & 0.810 $\pm$ 0.013 & 0.692 $\pm$ 0.018 & 0.629 $\pm$ 0.013 & 0.912 $\pm$ 0.037 & 0.841 $\pm$ 0.049 & 0.764 $\pm$ 0.048 & 0.752 $\pm$ 0.033 & 0.7912 & 5 \\
k4 & 0.924 $\pm$ 0.032 & 0.809 $\pm$ 0.015 & 0.694 $\pm$ 0.012 & 0.613 $\pm$ 0.032 & 0.885 $\pm$ 0.042 & 0.846 $\pm$ 0.046 & 0.770 $\pm$ 0.048 & 0.737 $\pm$ 0.035 & 0.7847 & 5 \\
k5 & 0.922 $\pm$ 0.033 & 0.806 $\pm$ 0.010 & 0.693 $\pm$ 0.012 & 0.597 $\pm$ 0.054 & 0.889 $\pm$ 0.021 & 0.834 $\pm$ 0.049 & 0.767 $\pm$ 0.059 & 0.754 $\pm$ 0.034 & 0.7829 & 4 \\
k6 & 0.932 $\pm$ 0.019 & 0.808 $\pm$ 0.018 & 0.685 $\pm$ 0.021 & 0.625 $\pm$ 0.015 & 0.865 $\pm$ 0.082 & 0.816 $\pm$ 0.036 & 0.763 $\pm$ 0.047 & 0.739 $\pm$ 0.026 & 0.7792 & 5 \\
k7 & 0.927 $\pm$ 0.029 & 0.804 $\pm$ 0.011 & 0.691 $\pm$ 0.016 & 0.619 $\pm$ 0.031 & 0.873 $\pm$ 0.056 & 0.823 $\pm$ 0.044 & 0.758 $\pm$ 0.055 & 0.738 $\pm$ 0.055 & 0.7789 & 5 \\
\bottomrule
\end{tabular}}
\end{table*}
\begin{table*}[t]
\centering
\caption{TPSA single-branch KAN differences relative to the matching all-MLP model. Values are AUC percentage points.}
\label{tab:si-tpsa-kan-deltas}
\resizebox{\textwidth}{!}{%
\begin{tabular}{lrrrrrrrrrrl}
\toprule
Branch & BBBP & Tox21 & ToxCast & SIDER & ClinTox & BACE & HIV & MUV & Macro & Wins & Assignment \\
\midrule
k0 & -0.16 & -0.11 & -0.91 & -0.22 & +2.11 & +3.00 & +0.64 & -1.37 & +0.372 & 3/8 & KAN \\
k1 & -0.88 & -0.22 & -0.52 & -0.51 & +0.01 & +2.70 & -0.10 & +0.18 & +0.082 & 3/8 & KAN \\
k2 & -1.17 & +0.09 & -0.43 & -0.02 & +3.10 & +0.93 & +0.44 & +1.69 & +0.578 & 5/8 & KAN \\
k3 & -0.93 & -0.53 & -0.61 & -0.31 & +1.90 & +0.38 & -0.10 & +2.13 & +0.240 & 3/8 & KAN \\
k4 & +0.47 & +0.01 & -0.59 & -0.41 & +0.31 & +0.29 & -1.90 & +0.79 & -0.131 & 5/8 & MLP \\
k5 & +0.25 & -0.25 & -0.65 & -0.29 & +3.41 & +1.81 & -0.53 & -2.96 & +0.099 & 3/8 & KAN \\
k6 & -0.90 & -0.28 & -0.21 & -0.54 & -2.10 & +0.53 & -0.58 & +0.07 & -0.500 & 2/8 & MLP \\
k7 & -0.54 & -0.84 & -0.50 & -0.11 & +1.98 & +0.85 & -1.10 & -0.47 & -0.092 & 2/8 & MLP \\
\bottomrule
\end{tabular}}
\end{table*}
\begin{table*}[t]
\centering
\caption{\MolLogP single-branch KAN differences relative to the matching all-MLP model. Values are AUC percentage points.}
\label{tab:si-logp-kan-deltas}
\resizebox{\textwidth}{!}{%
\begin{tabular}{lrrrrrrrrrrl}
\toprule
Branch & BBBP & Tox21 & ToxCast & SIDER & ClinTox & BACE & HIV & MUV & Macro & Wins & Assignment \\
\midrule
k0 & +0.00 & -0.84 & -0.75 & +0.63 & +0.40 & -1.75 & +0.56 & -2.62 & -0.546 & 4/8 & MLP \\
k1 & -0.68 & +0.05 & -0.11 & +2.14 & +3.56 & -0.73 & +1.23 & -0.24 & +0.651 & 4/8 & KAN \\
k2 & -0.16 & -0.10 & -0.90 & +0.84 & +3.01 & +0.14 & -1.15 & +0.43 & +0.263 & 4/8 & KAN \\
k3 & -0.60 & +0.12 & -0.61 & +2.30 & +4.42 & +0.17 & -0.58 & +1.26 & +0.810 & 5/8 & KAN \\
k4 & -1.22 & +0.01 & -0.39 & +0.66 & +1.76 & +0.69 & +0.04 & -0.25 & +0.163 & 5/8 & KAN \\
k5 & -1.37 & -0.27 & -0.51 & -0.94 & +2.19 & -0.51 & -0.22 & +1.47 & -0.021 & 2/8 & MLP \\
k6 & -0.40 & -0.07 & -1.28 & +1.89 & -0.27 & -2.32 & -0.59 & -0.04 & -0.385 & 1/8 & MLP \\
k7 & -0.94 & -0.55 & -0.70 & +1.25 & +0.57 & -1.65 & -1.12 & -0.21 & -0.420 & 2/8 & MLP \\
\bottomrule
\end{tabular}}
\end{table*}

\section{Fold-Level Legacy Eight-Layer Results}

Figures~\ref{fig:si-tpsa-folds} and \ref{fig:si-logp-folds} show the individual
scaffold-fold measurements. A horizontal segment marks the fold mean.
Tables~\ref{tab:si-tpsa-fold-results} and
\ref{tab:si-logp-fold-results} retain all five numerical fold values for every
legacy eight-layer-aligned model and dataset combination. These tables audit
the available files and do not imply complete split-row coverage.

\begin{figure*}[t]
\centering
\includegraphics[width=\textwidth]{figS07_tpsa_fold_distributions.pdf}
\caption{TPSA legacy eight-layer fold-level ROC-AUC values for the five main model
roles on each downstream dataset.}
\label{fig:si-tpsa-folds}
\end{figure*}

\begin{figure*}[t]
\centering
\includegraphics[width=\textwidth]{figS08_logp_fold_distributions.pdf}
\caption{MolLogP legacy eight-layer fold-level ROC-AUC values for the five main model
roles on each downstream dataset.}
\label{fig:si-logp-folds}
\end{figure*}

\begin{landscape}
\begin{longtable}{llrrrrrr}
\caption{TPSA legacy eight-layer-aligned fold-level results. Every model and dataset cell contains fold files 1 through 5 under the historical incomplete-batch evaluation protocol.}
\label{tab:si-tpsa-fold-results}\\
\toprule
Dataset & Role & Fold 1 & Fold 2 & Fold 3 & Fold 4 & Fold 5 & Mean $\pm$ SD \\
\midrule
\endfirsthead
\toprule
Dataset & Role & Fold 1 & Fold 2 & Fold 3 & Fold 4 & Fold 5 & Mean $\pm$ SD \\
\midrule
\endhead
BBBP & Base & 0.9259 & 0.9674 & 0.9362 & 0.9467 & 0.9659 & 0.9484 $\pm$ 0.0182 \\
 & I1 & 0.9274 & 0.9665 & 0.9254 & 0.9337 & 0.9698 & 0.9446 $\pm$ 0.0218 \\
 & I1+I2 & 0.9227 & 0.9754 & 0.9239 & 0.9327 & 0.9654 & 0.9440 $\pm$ 0.0246 \\
 & I1+I3 & 0.9362 & 0.9741 & 0.9183 & 0.9207 & 0.9686 & 0.9436 $\pm$ 0.0263 \\
 & Full & 0.9271 & 0.9652 & 0.9241 & 0.9271 & 0.9557 & 0.9399 $\pm$ 0.0192 \\
 & all-KAN & 0.9210 & 0.9752 & 0.8915 & 0.9316 & 0.9642 & 0.9367 $\pm$ 0.0338 \\
\midrule
Tox21 & Base & 0.7974 & 0.8288 & 0.7961 & 0.8410 & 0.8149 & 0.8156 $\pm$ 0.0196 \\
 & I1 & 0.7806 & 0.8279 & 0.7906 & 0.8204 & 0.7984 & 0.8036 $\pm$ 0.0200 \\
 & I1+I2 & 0.7958 & 0.8179 & 0.7896 & 0.8214 & 0.8242 & 0.8098 $\pm$ 0.0159 \\
 & I1+I3 & 0.8015 & 0.8287 & 0.8011 & 0.8089 & 0.8125 & 0.8105 $\pm$ 0.0113 \\
 & Full & 0.7809 & 0.8230 & 0.8063 & 0.8206 & 0.8188 & 0.8099 $\pm$ 0.0174 \\
 & all-KAN & 0.7909 & 0.8199 & 0.7887 & 0.8258 & 0.7671 & 0.7985 $\pm$ 0.0242 \\
\midrule
ToxCast & Base & 0.7059 & 0.7010 & 0.6709 & 0.7051 & 0.6952 & 0.6956 $\pm$ 0.0145 \\
 & I1 & 0.7190 & 0.7049 & 0.6695 & 0.7025 & 0.6663 & 0.6924 $\pm$ 0.0233 \\
 & I1+I2 & 0.7059 & 0.7012 & 0.6708 & 0.6780 & 0.6785 & 0.6869 $\pm$ 0.0156 \\
 & I1+I3 & 0.7081 & 0.7146 & 0.6826 & 0.6997 & 0.6897 & 0.6989 $\pm$ 0.0131 \\
 & Full & 0.6983 & 0.7178 & 0.6621 & 0.6910 & 0.6732 & 0.6885 $\pm$ 0.0217 \\
 & all-KAN & 0.7187 & 0.7157 & 0.6501 & 0.6923 & 0.6690 & 0.6892 $\pm$ 0.0297 \\
\midrule
SIDER & Base & 0.5909 & 0.6600 & 0.6387 & 0.6432 & 0.6223 & 0.6310 $\pm$ 0.0261 \\
 & I1 & 0.5813 & 0.6504 & 0.6409 & 0.6321 & 0.6125 & 0.6234 $\pm$ 0.0274 \\
 & I1+I2 & 0.5511 & 0.6517 & 0.6432 & 0.6142 & 0.6212 & 0.6163 $\pm$ 0.0395 \\
 & I1+I3 & 0.5787 & 0.6294 & 0.6815 & 0.6250 & 0.6222 & 0.6274 $\pm$ 0.0365 \\
 & Full & 0.5542 & 0.6424 & 0.6252 & 0.6170 & 0.6216 & 0.6121 $\pm$ 0.0337 \\
 & all-KAN & 0.5770 & 0.6491 & 0.6204 & 0.6334 & 0.5932 & 0.6146 $\pm$ 0.0294 \\
\midrule
ClinTox & Base & 0.7914 & 0.8885 & 0.9057 & 0.8423 & 0.8270 & 0.8510 $\pm$ 0.0464 \\
 & I1 & 0.8769 & 0.9384 & 0.9428 & 0.8978 & 0.8365 & 0.8985 $\pm$ 0.0444 \\
 & I1+I2 & 0.8838 & 0.8922 & 0.9202 & 0.8828 & 0.8292 & 0.8817 $\pm$ 0.0330 \\
 & I1+I3 & 0.8929 & 0.9439 & 0.8377 & 0.8548 & 0.8191 & 0.8697 $\pm$ 0.0496 \\
 & Full & 0.8297 & 0.8984 & 0.8918 & 0.8583 & 0.8298 & 0.8616 $\pm$ 0.0328 \\
 & all-KAN & 0.8424 & 0.8898 & 0.8919 & 0.9256 & 0.8844 & 0.8868 $\pm$ 0.0296 \\
\midrule
BACE & Base & 0.8843 & 0.8392 & 0.8940 & 0.8246 & 0.9479 & 0.8780 $\pm$ 0.0488 \\
 & I1 & 0.8721 & 0.8403 & 0.8769 & 0.8372 & 0.9283 & 0.8709 $\pm$ 0.0368 \\
 & I1+I2 & 0.8902 & 0.8344 & 0.8977 & 0.8246 & 0.9344 & 0.8763 $\pm$ 0.0460 \\
 & I1+I3 & 0.8848 & 0.8395 & 0.9050 & 0.8385 & 0.9217 & 0.8779 $\pm$ 0.0379 \\
 & Full & 0.8760 & 0.8062 & 0.9118 & 0.8051 & 0.8203 & 0.8439 $\pm$ 0.0478 \\
 & all-KAN & 0.8914 & 0.8197 & 0.8974 & 0.8367 & 0.9001 & 0.8691 $\pm$ 0.0379 \\
\midrule
HIV & Base & 0.7598 & 0.8423 & 0.7847 & 0.8116 & 0.7836 & 0.7964 $\pm$ 0.0315 \\
 & I1 & 0.7678 & 0.8008 & 0.7848 & 0.7670 & 0.7393 & 0.7719 $\pm$ 0.0229 \\
 & I1+I2 & 0.7720 & 0.8356 & 0.7857 & 0.7722 & 0.7591 & 0.7849 $\pm$ 0.0298 \\
 & I1+I3 & 0.7957 & 0.8636 & 0.7908 & 0.8168 & 0.7511 & 0.8036 $\pm$ 0.0411 \\
 & Full & 0.7450 & 0.8406 & 0.7864 & 0.7922 & 0.7617 & 0.7852 $\pm$ 0.0363 \\
 & all-KAN & 0.7816 & 0.8263 & 0.7756 & 0.7687 & 0.7466 & 0.7798 $\pm$ 0.0292 \\
\midrule
MUV & Base & 0.7802 & 0.7690 & 0.7101 & 0.8371 & 0.7638 & 0.7721 $\pm$ 0.0453 \\
 & I1 & 0.7020 & 0.7714 & 0.7093 & 0.8062 & 0.6810 & 0.7340 $\pm$ 0.0526 \\
 & I1+I2 & 0.7616 & 0.7185 & 0.7110 & 0.8027 & 0.7069 & 0.7401 $\pm$ 0.0412 \\
 & I1+I3 & 0.7093 & 0.7619 & 0.7877 & 0.7808 & 0.7342 & 0.7547 $\pm$ 0.0328 \\
 & Full & 0.6835 & 0.6808 & 0.7578 & 0.7270 & 0.6843 & 0.7067 $\pm$ 0.0344 \\
 & all-KAN & 0.7349 & 0.7566 & 0.7714 & 0.7520 & 0.6815 & 0.7393 $\pm$ 0.0349 \\
\bottomrule
\end{longtable}
\end{landscape}
\begin{landscape}
\begin{longtable}{llrrrrrr}
\caption{MolLogP legacy eight-layer-aligned fold-level results. Every model and dataset cell contains fold files 1 through 5 under the historical incomplete-batch evaluation protocol.}
\label{tab:si-logp-fold-results}\\
\toprule
Dataset & Role & Fold 1 & Fold 2 & Fold 3 & Fold 4 & Fold 5 & Mean $\pm$ SD \\
\midrule
\endfirsthead
\toprule
Dataset & Role & Fold 1 & Fold 2 & Fold 3 & Fold 4 & Fold 5 & Mean $\pm$ SD \\
\midrule
\endhead
BBBP & Base & 0.9259 & 0.9674 & 0.9362 & 0.9467 & 0.9659 & 0.9484 $\pm$ 0.0182 \\
 & I1 & 0.9357 & 0.9370 & 0.8900 & 0.9252 & 0.9604 & 0.9297 $\pm$ 0.0256 \\
 & I1+I2 & 0.9371 & 0.9421 & 0.9302 & 0.9303 & 0.9489 & 0.9377 $\pm$ 0.0080 \\
 & I1+I3 & 0.9148 & 0.9313 & 0.9229 & 0.9475 & 0.9591 & 0.9351 $\pm$ 0.0181 \\
 & Full & 0.9206 & 0.9747 & 0.9268 & 0.9564 & 0.9550 & 0.9467 $\pm$ 0.0225 \\
 & all-KAN & 0.9398 & 0.9506 & 0.8963 & 0.9127 & 0.9522 & 0.9303 $\pm$ 0.0247 \\
\midrule
Tox21 & Base & 0.7974 & 0.8288 & 0.7961 & 0.8410 & 0.8149 & 0.8156 $\pm$ 0.0196 \\
 & I1 & 0.7921 & 0.8186 & 0.7900 & 0.8225 & 0.8064 & 0.8059 $\pm$ 0.0148 \\
 & I1+I2 & 0.8017 & 0.8313 & 0.8008 & 0.8306 & 0.8084 & 0.8146 $\pm$ 0.0153 \\
 & I1+I3 & 0.7952 & 0.8158 & 0.7937 & 0.8227 & 0.8094 & 0.8074 $\pm$ 0.0127 \\
 & Full & 0.7811 & 0.8077 & 0.7938 & 0.8234 & 0.8172 & 0.8046 $\pm$ 0.0173 \\
 & all-KAN & 0.7900 & 0.8066 & 0.7745 & 0.8210 & 0.7873 & 0.7959 $\pm$ 0.0181 \\
\midrule
ToxCast & Base & 0.7059 & 0.7010 & 0.6709 & 0.7051 & 0.6952 & 0.6956 $\pm$ 0.0145 \\
 & I1 & 0.7145 & 0.6997 & 0.6650 & 0.7110 & 0.6679 & 0.6916 $\pm$ 0.0237 \\
 & I1+I2 & 0.7095 & 0.7037 & 0.6770 & 0.7028 & 0.6787 & 0.6943 $\pm$ 0.0153 \\
 & I1+I3 & 0.7071 & 0.6980 & 0.6539 & 0.6879 & 0.6757 & 0.6845 $\pm$ 0.0207 \\
 & Full & 0.7070 & 0.6969 & 0.6746 & 0.7005 & 0.6756 & 0.6909 $\pm$ 0.0149 \\
 & all-KAN & 0.7141 & 0.6924 & 0.6820 & 0.7014 & 0.6775 & 0.6935 $\pm$ 0.0148 \\
\midrule
SIDER & Base & 0.5909 & 0.6600 & 0.6387 & 0.6432 & 0.6223 & 0.6310 $\pm$ 0.0261 \\
 & I1 & 0.6062 & 0.6266 & 0.6452 & 0.6173 & 0.6213 & 0.6233 $\pm$ 0.0143 \\
 & I1+I2 & 0.5760 & 0.6051 & 0.6530 & 0.6245 & 0.6271 & 0.6171 $\pm$ 0.0286 \\
 & I1+I3 & 0.5740 & 0.6497 & 0.6520 & 0.5913 & 0.6018 & 0.6137 $\pm$ 0.0353 \\
 & Full & 0.5867 & 0.6295 & 0.6471 & 0.6103 & 0.6371 & 0.6221 $\pm$ 0.0239 \\
 & all-KAN & 0.6017 & 0.5884 & 0.6660 & 0.6238 & 0.5966 & 0.6153 $\pm$ 0.0312 \\
\midrule
ClinTox & Base & 0.7914 & 0.8885 & 0.9057 & 0.8423 & 0.8270 & 0.8510 $\pm$ 0.0464 \\
 & I1 & 0.8370 & 0.8903 & 0.8680 & 0.9082 & 0.8917 & 0.8790 $\pm$ 0.0275 \\
 & I1+I2 & 0.8218 & 0.8216 & 0.8228 & 0.8669 & 0.8452 & 0.8356 $\pm$ 0.0201 \\
 & I1+I3 & 0.9009 & 0.8076 & 0.8379 & 0.8926 & 0.8874 & 0.8653 $\pm$ 0.0406 \\
 & Full & 0.8578 & 0.8399 & 0.8149 & 0.8477 & 0.8646 & 0.8450 $\pm$ 0.0193 \\
 & all-KAN & 0.9145 & 0.9114 & 0.8802 & 0.8844 & 0.8819 & 0.8945 $\pm$ 0.0170 \\
\midrule
BACE & Base & 0.8843 & 0.8392 & 0.8940 & 0.8246 & 0.9479 & 0.8780 $\pm$ 0.0488 \\
 & I1 & 0.8355 & 0.8179 & 0.9053 & 0.8349 & 0.8466 & 0.8480 $\pm$ 0.0336 \\
 & I1+I2 & 0.8412 & 0.8159 & 0.9033 & 0.8564 & 0.9079 & 0.8649 $\pm$ 0.0399 \\
 & I1+I3 & 0.8453 & 0.8197 & 0.8757 & 0.8559 & 0.8759 & 0.8545 $\pm$ 0.0234 \\
 & Full & 0.8431 & 0.8028 & 0.8713 & 0.8459 & 0.8807 & 0.8488 $\pm$ 0.0303 \\
 & all-KAN & 0.7985 & 0.8064 & 0.9270 & 0.8269 & 0.8965 & 0.8511 $\pm$ 0.0574 \\
\midrule
HIV & Base & 0.7598 & 0.8423 & 0.7847 & 0.8116 & 0.7836 & 0.7964 $\pm$ 0.0315 \\
 & I1 & 0.7738 & 0.8409 & 0.7945 & 0.7490 & 0.7499 & 0.7816 $\pm$ 0.0381 \\
 & I1+I2 & 0.7341 & 0.8572 & 0.8042 & 0.7640 & 0.7450 & 0.7809 $\pm$ 0.0503 \\
 & I1+I3 & 0.7640 & 0.8432 & 0.8136 & 0.7731 & 0.7493 & 0.7887 $\pm$ 0.0387 \\
 & Full & 0.7215 & 0.8525 & 0.7951 & 0.7822 & 0.7534 & 0.7809 $\pm$ 0.0490 \\
 & all-KAN & 0.7318 & 0.8241 & 0.7715 & 0.7463 & 0.7279 & 0.7603 $\pm$ 0.0396 \\
\midrule
MUV & Base & 0.7802 & 0.7690 & 0.7101 & 0.8371 & 0.7638 & 0.7721 $\pm$ 0.0453 \\
 & I1 & 0.7154 & 0.7448 & 0.7530 & 0.7555 & 0.7173 & 0.7372 $\pm$ 0.0195 \\
 & I1+I2 & 0.7342 & 0.7514 & 0.7555 & 0.7296 & 0.7797 & 0.7501 $\pm$ 0.0199 \\
 & I1+I3 & 0.7293 & 0.7717 & 0.7355 & 0.7689 & 0.7311 & 0.7473 $\pm$ 0.0212 \\
 & Full & 0.6773 & 0.7196 & 0.7755 & 0.6845 & 0.7406 & 0.7195 $\pm$ 0.0406 \\
 & all-KAN & 0.7039 & 0.7807 & 0.7220 & 0.6662 & 0.7179 & 0.7182 $\pm$ 0.0413 \\
\bottomrule
\end{longtable}
\end{landscape}

\section{Per-Dataset Conditional Ablation and Statistical Tests}

Tables~\ref{tab:si-tpsa-conditional} and
\ref{tab:si-logp-conditional} report every dataset-level component contrast.
The Wilcoxon analysis in Table~\ref{tab:si-statistics} is secondary to the
effect sizes and hierarchical bootstrap intervals in the main text because
only eight datasets are available.

\begin{table*}[t]
\centering
\caption{TPSA per-dataset conditional ablation differences. Values are matched AUC percentage points.}
\label{tab:si-tpsa-conditional}
\resizebox{\textwidth}{!}{%
\begin{tabular}{lrrrrrrrrrr}
\toprule
Contrast & BBBP & Tox21 & ToxCast & SIDER & ClinTox & BACE & HIV & MUV & Macro & Wins \\
\midrule
$\Delta I_1$ & -0.38 & -1.20 & -0.32 & -0.76 & +4.75 & -0.71 & -2.45 & -3.81 & -0.609 & 1/8 \\
$\Delta I_2$ & -0.05 & +0.62 & -0.55 & -0.72 & -1.68 & +0.53 & +1.30 & +0.61 & +0.007 & 4/8 \\
$\Delta I_3$ & -0.10 & +0.69 & +0.65 & +0.40 & -2.88 & +0.69 & +3.16 & +2.08 & +0.587 & 6/8 \\
Full vs Base & -0.86 & -0.57 & -0.71 & -1.89 & +1.06 & -3.41 & -1.12 & -6.54 & -1.756 & 1/8 \\
$\Delta_{\mathrm{interaction}}$ & -0.32 & -0.68 & -0.49 & -0.82 & +0.87 & -3.93 & -3.14 & -5.42 & -1.740 & \\
\bottomrule
\end{tabular}}
\end{table*}
\begin{table*}[t]
\centering
\caption{MolLogP per-dataset conditional ablation differences. Values are matched AUC percentage points.}
\label{tab:si-logp-conditional}
\resizebox{\textwidth}{!}{%
\begin{tabular}{lrrrrrrrrrr}
\toprule
Contrast & BBBP & Tox21 & ToxCast & SIDER & ClinTox & BACE & HIV & MUV & Macro & Wins \\
\midrule
$\Delta I_1$ & -1.88 & -0.97 & -0.40 & -0.77 & +2.81 & -3.00 & -1.48 & -3.49 & -1.146 & 1/8 \\
$\Delta I_2$ & +0.80 & +0.86 & +0.27 & -0.62 & -4.34 & +1.69 & -0.07 & +1.29 & -0.013 & 5/8 \\
$\Delta I_3$ & +0.55 & +0.14 & -0.71 & -0.96 & -1.37 & +0.65 & +0.70 & +1.01 & +0.001 & 5/8 \\
Full vs Base & -0.17 & -1.10 & -0.47 & -0.89 & -0.60 & -2.92 & -1.55 & -5.26 & -1.619 & 0/8 \\
$\Delta_{\mathrm{interaction}}$ & +0.35 & -1.14 & +0.37 & +1.46 & +2.31 & -2.26 & -0.70 & -4.07 & -0.461 & \\
\bottomrule
\end{tabular}}
\end{table*}
\begin{table}[t]
\centering
\caption{Two-sided Wilcoxon signed-rank tests on complete shared-dataset mean differences. Holm adjustment was applied across all ten comparisons.}
\label{tab:si-statistics}
\begin{tabular}{llrrrr}
\toprule
Descriptor & Comparison & Data sets & Statistic & $p$ & Holm $p$ \\
\midrule
TPSA & I1 vs Base & 8 & 8.0 & 0.1953 & 1.0000 \\
TPSA & I2 conditional effect & 8 & 18.0 & 1.0000 & 1.0000 \\
TPSA & I3 conditional effect & 8 & 8.0 & 0.1953 & 1.0000 \\
TPSA & Full vs Base & 8 & 4.0 & 0.0547 & 0.4922 \\
TPSA & selected-KAN vs all-KAN & 8 & 8.0 & 0.1953 & 1.0000 \\
MolLogP & I1 vs Base & 8 & 6.0 & 0.1094 & 0.8750 \\
MolLogP & I2 conditional effect & 8 & 12.0 & 0.4609 & 1.0000 \\
MolLogP & I3 conditional effect & 8 & 17.0 & 0.9453 & 1.0000 \\
MolLogP & Full vs Base & 8 & 0.0 & 0.0078 & 0.0781 \\
MolLogP & selected-KAN vs all-KAN & 8 & 8.0 & 0.1953 & 1.0000 \\
\bottomrule
\end{tabular}
\end{table}

\section{Auxiliary Controls}

The N4 Base is a depth-sensitivity reference. Routed top1 to top4 models use
learned latent routing at all Transformer layers and are therefore not
layer-matched causal controls for the last-layer descriptor assignment.
Incomplete five-fold cells remain blank.

\begin{figure*}[t]
\centering
\includegraphics[width=\textwidth]{figS09_auxiliary_controls.pdf}
\caption{Auxiliary depth, learned-routing, and all-KAN controls under their
recorded protocols. Blank cells contain fewer than five folds.}
\label{fig:si-controls}
\end{figure*}

\begin{table*}[t]
\centering
\caption{Auxiliary depth and learned-routing controls. Empty cells contain fewer than five completed folds. Routed models replace FFNs across all layers and are not layer-matched causal controls for MoFE-last.}
\label{tab:si-controls}
\resizebox{\textwidth}{!}{%
\begin{tabular}{lrrrrrrrr}
\toprule
Control & BBBP & Tox21 & ToxCast & SIDER & ClinTox & BACE & HIV & MUV \\
\midrule
N4 Base & 0.947 $\pm$ 0.018 & 0.818 $\pm$ 0.020 & 0.699 $\pm$ 0.019 & 0.614 $\pm$ 0.020 & 0.873 $\pm$ 0.057 & 0.876 $\pm$ 0.041 & 0.802 $\pm$ 0.023 & 0.752 $\pm$ 0.024 \\
Routed top1 & 0.949 $\pm$ 0.020 & 0.805 $\pm$ 0.015 & 0.697 $\pm$ 0.017 & 0.625 $\pm$ 0.028 & 0.882 $\pm$ 0.054 & 0.875 $\pm$ 0.038 & \pendingresult & 0.755 $\pm$ 0.038 \\
Routed top2 & 0.941 $\pm$ 0.023 & 0.813 $\pm$ 0.014 & 0.690 $\pm$ 0.014 & 0.634 $\pm$ 0.029 & 0.873 $\pm$ 0.043 & 0.866 $\pm$ 0.053 & 0.804 $\pm$ 0.030 & 0.760 $\pm$ 0.028 \\
Routed top3 & 0.944 $\pm$ 0.032 & 0.809 $\pm$ 0.018 & \pendingresult & 0.628 $\pm$ 0.023 & 0.880 $\pm$ 0.065 & 0.877 $\pm$ 0.043 & \pendingresult & 0.774 $\pm$ 0.037 \\
Routed top4 & 0.939 $\pm$ 0.029 & \pendingresult & \pendingresult & 0.631 $\pm$ 0.014 & 0.876 $\pm$ 0.047 & 0.863 $\pm$ 0.052 & 0.808 $\pm$ 0.035 & 0.766 $\pm$ 0.042 \\
\bottomrule
\end{tabular}}
\end{table*}

\section{Complete Embedding Separation Metrics}

Quantitative metrics are calculated in model-specific PCA spaces. Means and
sample SDs require all five folds. Incomplete five-fold entries remain blank.

\begin{landscape}
\begin{longtable}{lllrrrrrr}
\caption{Descriptor-region separation metrics across all available embedding folds. Means and sample SDs are shown only for complete five-fold entries.}
\label{tab:si-embedding-all}\\
\toprule
Partition & Role & Dataset & Folds & Silhouette & DB index & Intra & Inter & Ratio \\
\midrule
\endfirsthead
\toprule
Partition & Role & Dataset & Folds & Silhouette & DB index & Intra & Inter & Ratio \\
\midrule
\endhead
TPSA & Base & BBBP & 5 & -0.035 $\pm$ 0.031 & 4.929 $\pm$ 0.480 & 1.252 $\pm$ 0.046 & 1.413 $\pm$ 0.015 & 1.130 $\pm$ 0.038 \\
TPSA & Base & Tox21 & 5 & -0.015 $\pm$ 0.008 & 8.063 $\pm$ 0.708 & 1.337 $\pm$ 0.013 & 1.410 $\pm$ 0.002 & 1.055 $\pm$ 0.011 \\
TPSA & Base & ToxCast & 5 & -0.014 $\pm$ 0.007 & 8.469 $\pm$ 1.052 & 1.326 $\pm$ 0.010 & 1.409 $\pm$ 0.003 & 1.062 $\pm$ 0.008 \\
TPSA & Base & SIDER & 5 & -0.016 $\pm$ 0.014 & 4.758 $\pm$ 0.126 & 1.324 $\pm$ 0.026 & 1.421 $\pm$ 0.005 & 1.073 $\pm$ 0.024 \\
TPSA & Base & ClinTox & 5 & -0.032 $\pm$ 0.013 & 4.738 $\pm$ 0.587 & 1.309 $\pm$ 0.029 & 1.415 $\pm$ 0.004 & 1.081 $\pm$ 0.021 \\
TPSA & Base & BACE & 5 & -0.000 $\pm$ 0.020 & 3.351 $\pm$ 0.541 & 1.259 $\pm$ 0.025 & 1.432 $\pm$ 0.010 & 1.138 $\pm$ 0.030 \\
TPSA & Base & HIV & 5 & -0.017 $\pm$ 0.003 & 10.853 $\pm$ 0.317 & 1.356 $\pm$ 0.008 & 1.410 $\pm$ 0.002 & 1.040 $\pm$ 0.006 \\
TPSA & Base & MUV & 5 & -0.058 $\pm$ 0.012 & 8.843 $\pm$ 0.691 & 1.388 $\pm$ 0.005 & 1.410 $\pm$ 0.001 & 1.016 $\pm$ 0.004 \\
TPSA & I1 & BBBP & 5 & 0.035 $\pm$ 0.019 & 3.137 $\pm$ 0.299 & 1.195 $\pm$ 0.036 & 1.433 $\pm$ 0.005 & 1.200 $\pm$ 0.039 \\
TPSA & I1 & Tox21 & 5 & 0.031 $\pm$ 0.008 & 4.238 $\pm$ 0.212 & 1.295 $\pm$ 0.015 & 1.419 $\pm$ 0.003 & 1.096 $\pm$ 0.015 \\
TPSA & I1 & ToxCast & 5 & 0.021 $\pm$ 0.008 & 4.531 $\pm$ 0.309 & 1.288 $\pm$ 0.013 & 1.416 $\pm$ 0.003 & 1.099 $\pm$ 0.013 \\
TPSA & I1 & SIDER & 5 & 0.035 $\pm$ 0.011 & 3.373 $\pm$ 0.317 & 1.296 $\pm$ 0.029 & 1.428 $\pm$ 0.003 & 1.102 $\pm$ 0.026 \\
TPSA & I1 & ClinTox & 5 & 0.013 $\pm$ 0.030 & 3.120 $\pm$ 0.243 & 1.265 $\pm$ 0.032 & 1.423 $\pm$ 0.005 & 1.125 $\pm$ 0.029 \\
TPSA & I1 & BACE & 5 & 0.083 $\pm$ 0.029 & 2.306 $\pm$ 0.242 & 1.210 $\pm$ 0.034 & 1.454 $\pm$ 0.013 & 1.203 $\pm$ 0.043 \\
TPSA & I1 & HIV & 5 & 0.020 $\pm$ 0.014 & 4.903 $\pm$ 0.842 & 1.317 $\pm$ 0.024 & 1.418 $\pm$ 0.003 & 1.077 $\pm$ 0.022 \\
TPSA & I1 & MUV & 5 & -0.019 $\pm$ 0.011 & 3.922 $\pm$ 0.329 & 1.354 $\pm$ 0.006 & 1.422 $\pm$ 0.002 & 1.050 $\pm$ 0.006 \\
TPSA & I1+I2 & BBBP & 5 & 0.064 $\pm$ 0.032 & 2.500 $\pm$ 0.106 & 1.131 $\pm$ 0.049 & 1.433 $\pm$ 0.007 & 1.269 $\pm$ 0.054 \\
TPSA & I1+I2 & Tox21 & 5 & 0.051 $\pm$ 0.011 & 3.299 $\pm$ 0.168 & 1.262 $\pm$ 0.019 & 1.424 $\pm$ 0.002 & 1.129 $\pm$ 0.018 \\
TPSA & I1+I2 & ToxCast & 5 & 0.052 $\pm$ 0.007 & 3.381 $\pm$ 0.257 & 1.255 $\pm$ 0.014 & 1.421 $\pm$ 0.004 & 1.133 $\pm$ 0.014 \\
TPSA & I1+I2 & SIDER & 5 & 0.052 $\pm$ 0.028 & 3.036 $\pm$ 0.503 & 1.275 $\pm$ 0.051 & 1.431 $\pm$ 0.006 & 1.124 $\pm$ 0.050 \\
TPSA & I1+I2 & ClinTox & 5 & 0.045 $\pm$ 0.043 & 2.756 $\pm$ 0.275 & 1.228 $\pm$ 0.032 & 1.428 $\pm$ 0.006 & 1.164 $\pm$ 0.034 \\
TPSA & I1+I2 & BACE & 5 & 0.105 $\pm$ 0.036 & 2.211 $\pm$ 0.449 & 1.183 $\pm$ 0.052 & 1.461 $\pm$ 0.018 & 1.237 $\pm$ 0.069 \\
TPSA & I1+I2 & HIV & 5 & 0.037 $\pm$ 0.016 & 3.824 $\pm$ 0.426 & 1.287 $\pm$ 0.021 & 1.421 $\pm$ 0.002 & 1.105 $\pm$ 0.019 \\
TPSA & I1+I2 & MUV & 5 & -0.003 $\pm$ 0.016 & 3.180 $\pm$ 0.224 & 1.324 $\pm$ 0.011 & 1.431 $\pm$ 0.002 & 1.080 $\pm$ 0.011 \\
TPSA & I1+I3 & BBBP & 5 & 0.286 $\pm$ 0.108 & 1.302 $\pm$ 0.339 & 0.935 $\pm$ 0.114 & 1.476 $\pm$ 0.020 & 1.600 $\pm$ 0.216 \\
TPSA & I1+I3 & Tox21 & 5 & 0.298 $\pm$ 0.105 & 1.486 $\pm$ 0.432 & 0.985 $\pm$ 0.136 & 1.472 $\pm$ 0.019 & 1.518 $\pm$ 0.213 \\
TPSA & I1+I3 & ToxCast & 5 & 0.074 $\pm$ 0.019 & 2.846 $\pm$ 0.306 & 1.225 $\pm$ 0.029 & 1.424 $\pm$ 0.005 & 1.163 $\pm$ 0.029 \\
TPSA & I1+I3 & SIDER & 5 & 0.242 $\pm$ 0.142 & 1.748 $\pm$ 0.748 & 1.042 $\pm$ 0.185 & 1.465 $\pm$ 0.024 & 1.445 $\pm$ 0.278 \\
TPSA & I1+I3 & ClinTox & 5 & 0.175 $\pm$ 0.164 & 1.889 $\pm$ 0.625 & 1.084 $\pm$ 0.193 & 1.453 $\pm$ 0.026 & 1.388 $\pm$ 0.334 \\
TPSA & I1+I3 & BACE & 5 & 0.406 $\pm$ 0.181 & 0.991 $\pm$ 0.450 & 0.842 $\pm$ 0.222 & 1.544 $\pm$ 0.057 & 1.957 $\pm$ 0.591 \\
TPSA & I1+I3 & HIV & 5 & 0.398 $\pm$ 0.144 & 1.217 $\pm$ 0.384 & 0.821 $\pm$ 0.198 & 1.487 $\pm$ 0.020 & 1.900 $\pm$ 0.468 \\
TPSA & I1+I3 & MUV & 5 & 0.090 $\pm$ 0.051 & 2.096 $\pm$ 0.461 & 1.229 $\pm$ 0.064 & 1.458 $\pm$ 0.016 & 1.189 $\pm$ 0.074 \\
TPSA & Full & BBBP & 5 & 0.285 $\pm$ 0.149 & 1.393 $\pm$ 0.486 & 0.908 $\pm$ 0.141 & 1.471 $\pm$ 0.024 & 1.663 $\pm$ 0.337 \\
TPSA & Full & Tox21 & 5 & 0.392 $\pm$ 0.043 & 1.180 $\pm$ 0.129 & 0.858 $\pm$ 0.066 & 1.488 $\pm$ 0.007 & 1.744 $\pm$ 0.145 \\
TPSA & Full & ToxCast & 5 & 0.110 $\pm$ 0.029 & 2.510 $\pm$ 0.399 & 1.198 $\pm$ 0.038 & 1.434 $\pm$ 0.006 & 1.199 $\pm$ 0.042 \\
TPSA & Full & SIDER & 5 & 0.202 $\pm$ 0.104 & 1.754 $\pm$ 0.494 & 1.104 $\pm$ 0.142 & 1.458 $\pm$ 0.016 & 1.341 $\pm$ 0.206 \\
TPSA & Full & ClinTox & 5 & 0.245 $\pm$ 0.094 & 1.473 $\pm$ 0.346 & 0.999 $\pm$ 0.102 & 1.466 $\pm$ 0.017 & 1.480 $\pm$ 0.153 \\
TPSA & Full & BACE & 5 & 0.484 $\pm$ 0.165 & 0.845 $\pm$ 0.350 & 0.736 $\pm$ 0.225 & 1.561 $\pm$ 0.049 & 2.299 $\pm$ 0.740 \\
TPSA & Full & HIV & 5 & 0.444 $\pm$ 0.189 & 1.156 $\pm$ 0.635 & 0.748 $\pm$ 0.246 & 1.492 $\pm$ 0.030 & 2.144 $\pm$ 0.571 \\
TPSA & Full & MUV & 5 & 0.103 $\pm$ 0.098 & 2.026 $\pm$ 0.530 & 1.213 $\pm$ 0.107 & 1.462 $\pm$ 0.026 & 1.215 $\pm$ 0.139 \\
\MolLogP & Base & BBBP & 5 & -0.062 $\pm$ 0.020 & 5.920 $\pm$ 0.425 & 1.312 $\pm$ 0.029 & 1.401 $\pm$ 0.020 & 1.068 $\pm$ 0.025 \\
\MolLogP & Base & Tox21 & 5 & -0.026 $\pm$ 0.005 & 7.598 $\pm$ 0.508 & 1.338 $\pm$ 0.010 & 1.409 $\pm$ 0.002 & 1.053 $\pm$ 0.009 \\
\MolLogP & Base & ToxCast & 5 & -0.029 $\pm$ 0.007 & 7.116 $\pm$ 0.646 & 1.310 $\pm$ 0.028 & 1.411 $\pm$ 0.002 & 1.077 $\pm$ 0.024 \\
\MolLogP & Base & SIDER & 5 & -0.018 $\pm$ 0.013 & 4.946 $\pm$ 0.261 & 1.343 $\pm$ 0.010 & 1.417 $\pm$ 0.004 & 1.055 $\pm$ 0.011 \\
\MolLogP & Base & ClinTox & 5 & -0.042 $\pm$ 0.017 & 5.296 $\pm$ 0.145 & 1.343 $\pm$ 0.008 & 1.408 $\pm$ 0.008 & 1.049 $\pm$ 0.005 \\
\MolLogP & Base & BACE & 5 & -0.050 $\pm$ 0.040 & 4.256 $\pm$ 0.566 & 1.338 $\pm$ 0.032 & 1.404 $\pm$ 0.008 & 1.050 $\pm$ 0.031 \\
\MolLogP & Base & HIV & 5 & -0.021 $\pm$ 0.007 & 9.854 $\pm$ 0.737 & 1.349 $\pm$ 0.014 & 1.412 $\pm$ 0.001 & 1.047 $\pm$ 0.011 \\
\MolLogP & Base & MUV & 5 & -0.059 $\pm$ 0.011 & 8.738 $\pm$ 0.855 & 1.389 $\pm$ 0.004 & 1.410 $\pm$ 0.001 & 1.015 $\pm$ 0.003 \\
\MolLogP & I1 & BBBP & 5 & 0.014 $\pm$ 0.060 & 3.376 $\pm$ 0.441 & 1.240 $\pm$ 0.040 & 1.420 $\pm$ 0.022 & 1.146 $\pm$ 0.048 \\
\MolLogP & I1 & Tox21 & 5 & 0.045 $\pm$ 0.014 & 3.856 $\pm$ 0.307 & 1.262 $\pm$ 0.032 & 1.426 $\pm$ 0.004 & 1.130 $\pm$ 0.033 \\
\MolLogP & I1 & ToxCast & 5 & 0.016 $\pm$ 0.006 & 4.184 $\pm$ 0.196 & 1.257 $\pm$ 0.019 & 1.421 $\pm$ 0.001 & 1.130 $\pm$ 0.018 \\
\MolLogP & I1 & SIDER & 5 & 0.036 $\pm$ 0.027 & 3.347 $\pm$ 0.445 & 1.291 $\pm$ 0.056 & 1.429 $\pm$ 0.005 & 1.109 $\pm$ 0.053 \\
\MolLogP & I1 & ClinTox & 5 & 0.036 $\pm$ 0.029 & 3.206 $\pm$ 0.202 & 1.270 $\pm$ 0.043 & 1.423 $\pm$ 0.006 & 1.121 $\pm$ 0.042 \\
\MolLogP & I1 & BACE & 5 & 0.074 $\pm$ 0.044 & 2.357 $\pm$ 0.121 & 1.236 $\pm$ 0.055 & 1.444 $\pm$ 0.006 & 1.170 $\pm$ 0.058 \\
\MolLogP & I1 & HIV & 5 & 0.041 $\pm$ 0.017 & 4.585 $\pm$ 0.484 & 1.277 $\pm$ 0.023 & 1.427 $\pm$ 0.003 & 1.118 $\pm$ 0.023 \\
\MolLogP & I1 & MUV & 5 & 0.009 $\pm$ 0.040 & 3.444 $\pm$ 0.527 & 1.327 $\pm$ 0.024 & 1.431 $\pm$ 0.007 & 1.078 $\pm$ 0.024 \\
\MolLogP & I1+I2 & BBBP & 5 & 0.056 $\pm$ 0.022 & 2.608 $\pm$ 0.136 & 1.208 $\pm$ 0.040 & 1.428 $\pm$ 0.007 & 1.184 $\pm$ 0.039 \\
\MolLogP & I1+I2 & Tox21 & 5 & 0.061 $\pm$ 0.009 & 3.476 $\pm$ 0.293 & 1.238 $\pm$ 0.030 & 1.428 $\pm$ 0.004 & 1.154 $\pm$ 0.031 \\
\MolLogP & I1+I2 & ToxCast & 5 & 0.025 $\pm$ 0.005 & 3.589 $\pm$ 0.158 & 1.240 $\pm$ 0.031 & 1.422 $\pm$ 0.003 & 1.147 $\pm$ 0.031 \\
\MolLogP & I1+I2 & SIDER & 5 & 0.066 $\pm$ 0.045 & 2.841 $\pm$ 0.472 & 1.240 $\pm$ 0.083 & 1.433 $\pm$ 0.008 & 1.161 $\pm$ 0.087 \\
\MolLogP & I1+I2 & ClinTox & 5 & 0.079 $\pm$ 0.027 & 2.588 $\pm$ 0.205 & 1.221 $\pm$ 0.039 & 1.433 $\pm$ 0.004 & 1.174 $\pm$ 0.040 \\
\MolLogP & I1+I2 & BACE & 5 & 0.089 $\pm$ 0.073 & 2.213 $\pm$ 0.308 & 1.222 $\pm$ 0.091 & 1.448 $\pm$ 0.020 & 1.191 $\pm$ 0.104 \\
\MolLogP & I1+I2 & HIV & 5 & 0.072 $\pm$ 0.028 & 3.746 $\pm$ 0.379 & 1.222 $\pm$ 0.046 & 1.433 $\pm$ 0.007 & 1.174 $\pm$ 0.048 \\
\MolLogP & I1+I2 & MUV & 5 & 0.020 $\pm$ 0.022 & 3.064 $\pm$ 0.504 & 1.312 $\pm$ 0.021 & 1.437 $\pm$ 0.004 & 1.095 $\pm$ 0.021 \\
\MolLogP & I1+I3 & BBBP & 5 & 0.223 $\pm$ 0.170 & 1.693 $\pm$ 0.684 & 1.027 $\pm$ 0.158 & 1.455 $\pm$ 0.032 & 1.447 $\pm$ 0.245 \\
\MolLogP & I1+I3 & Tox21 & 5 & 0.310 $\pm$ 0.047 & 1.384 $\pm$ 0.167 & 0.943 $\pm$ 0.064 & 1.474 $\pm$ 0.006 & 1.570 $\pm$ 0.110 \\
\MolLogP & I1+I3 & ToxCast & 5 & 0.066 $\pm$ 0.010 & 2.856 $\pm$ 0.224 & 1.199 $\pm$ 0.025 & 1.431 $\pm$ 0.003 & 1.194 $\pm$ 0.027 \\
\MolLogP & I1+I3 & SIDER & 5 & 0.164 $\pm$ 0.116 & 2.120 $\pm$ 0.676 & 1.151 $\pm$ 0.149 & 1.449 $\pm$ 0.018 & 1.280 $\pm$ 0.198 \\
\MolLogP & I1+I3 & ClinTox & 5 & 0.281 $\pm$ 0.137 & 1.515 $\pm$ 0.495 & 0.976 $\pm$ 0.171 & 1.468 $\pm$ 0.020 & 1.541 $\pm$ 0.269 \\
\MolLogP & I1+I3 & BACE & 5 & 0.415 $\pm$ 0.176 & 1.006 $\pm$ 0.382 & 0.810 $\pm$ 0.250 & 1.549 $\pm$ 0.061 & 2.069 $\pm$ 0.650 \\
\MolLogP & I1+I3 & HIV & 5 & 0.341 $\pm$ 0.122 & 1.399 $\pm$ 0.452 & 0.915 $\pm$ 0.152 & 1.479 $\pm$ 0.019 & 1.651 $\pm$ 0.257 \\
\MolLogP & I1+I3 & MUV & 5 & 0.054 $\pm$ 0.061 & 2.492 $\pm$ 0.709 & 1.265 $\pm$ 0.063 & 1.451 $\pm$ 0.018 & 1.150 $\pm$ 0.072 \\
\MolLogP & Full & BBBP & 5 & 0.263 $\pm$ 0.069 & 1.440 $\pm$ 0.259 & 0.997 $\pm$ 0.074 & 1.465 $\pm$ 0.014 & 1.476 $\pm$ 0.119 \\
\MolLogP & Full & Tox21 & 5 & 0.322 $\pm$ 0.056 & 1.344 $\pm$ 0.161 & 0.932 $\pm$ 0.071 & 1.476 $\pm$ 0.007 & 1.592 $\pm$ 0.131 \\
\MolLogP & Full & ToxCast & 5 & 0.105 $\pm$ 0.013 & 2.443 $\pm$ 0.109 & 1.145 $\pm$ 0.016 & 1.439 $\pm$ 0.002 & 1.257 $\pm$ 0.018 \\
\MolLogP & Full & SIDER & 5 & 0.221 $\pm$ 0.147 & 1.900 $\pm$ 0.784 & 1.070 $\pm$ 0.196 & 1.457 $\pm$ 0.022 & 1.407 $\pm$ 0.312 \\
\MolLogP & Full & ClinTox & 5 & 0.279 $\pm$ 0.131 & 1.498 $\pm$ 0.413 & 0.957 $\pm$ 0.155 & 1.468 $\pm$ 0.018 & 1.567 $\pm$ 0.259 \\
\MolLogP & Full & BACE & 5 & 0.299 $\pm$ 0.138 & 1.258 $\pm$ 0.354 & 0.975 $\pm$ 0.191 & 1.516 $\pm$ 0.048 & 1.612 $\pm$ 0.368 \\
\MolLogP & Full & HIV & 5 & 0.446 $\pm$ 0.044 & 1.075 $\pm$ 0.078 & 0.772 $\pm$ 0.060 & 1.494 $\pm$ 0.005 & 1.946 $\pm$ 0.153 \\
\MolLogP & Full & MUV & 5 & 0.107 $\pm$ 0.093 & 1.974 $\pm$ 0.531 & 1.199 $\pm$ 0.087 & 1.470 $\pm$ 0.023 & 1.232 $\pm$ 0.112 \\
\bottomrule
\end{longtable}
\end{landscape}

\section{Pretraining Dynamics and Active Runtime}

Only the logged terminal-batch composite loss and seconds per epoch were
recorded consistently. The logged loss is the final mini-batch value from each
epoch and is not an epoch-average loss. The active-time curves exclude the
first epoch of each Slurm segment and values above the segment-specific
\(Q_3+3\,\mathrm{IQR}\) upper outer fence. Queue time between segments is
absent from the in-process timer.

\begin{figure*}[t]
\centering
\includegraphics[width=\textwidth]{figS10_pretraining_dynamics_tpsa.pdf}
\caption{Recovered TPSA terminal-batch composite-loss and active epoch-time
traces. Time curves
show 11-epoch rolling medians and interquartile bands after deterministic
timing exclusions.}
\label{fig:si-tpsa-training}
\end{figure*}

\begin{figure*}[t]
\centering
\includegraphics[width=\textwidth]{figS11_pretraining_dynamics_logp.pdf}
\caption{Recovered MolLogP terminal-batch composite-loss and active epoch-time
traces. Separate
graph reconstruction, SMILES MLM, and ConLoss epoch series were not available
in the stored logs.}
\label{fig:si-logp-training}
\end{figure*}

\begin{table*}[t]
\centering
\caption{Recovered active pretraining time. The first epoch of each Slurm segment and high-duration values above the segment-specific $Q_3+3\,\mathrm{IQR}$ outer fence are excluded. Scheduler wait between segments is not included. Projected active hours equal the retained median multiplied by 300.}
\label{tab:si-pretrain-runtime}
\resizebox{\textwidth}{!}{%
\begin{tabular}{lrrrrrr}
\toprule
Model & Raw epochs & Retained & Excluded & Median min/epoch & Active IQR & Projected active h \\
\midrule
Base & 300 & 293 & 7 & 15.15 & 15.10 to 15.21 & 75.7 \\
TPSA I1 & 300 & 272 & 28 & 15.31 & 15.06 to 15.34 & 76.6 \\
MolLogP I1 & 300 & 292 & 8 & 15.86 & 15.81 to 15.92 & 79.3 \\
TPSA I1+I2 & 300 & 299 & 1 & 17.20 & 17.05 to 17.37 & 86.0 \\
MolLogP I1+I2 & 300 & 299 & 1 & 17.21 & 16.97 to 17.40 & 86.0 \\
TPSA I1+I3 & 300 & 287 & 13 & 16.02 & 15.99 to 16.05 & 80.1 \\
MolLogP I1+I3 & 300 & 290 & 10 & 16.05 & 15.98 to 16.08 & 80.3 \\
TPSA Full & 300 & 299 & 1 & 17.20 & 17.05 to 17.37 & 86.0 \\
MolLogP Full & 300 & 299 & 1 & 16.78 & 16.46 to 17.28 & 83.9 \\
TPSA all-KAN & 300 & 290 & 10 & 17.53 & 17.43 to 17.82 & 87.7 \\
MolLogP all-KAN & 300 & 276 & 24 & 17.73 & 17.65 to 18.02 & 88.7 \\
\bottomrule
\end{tabular}}
\end{table*}

\section{Downstream Runtime and Completion}

Runtime records describe consumed Slurm allocations rather than a controlled
architecture benchmark. Retries and cancelled attempts are retained
separately in the table.

\begin{figure*}[t]
\centering
\includegraphics[width=\textwidth]{figS12_finetuning_runtime.pdf}
\caption{Observed completed legacy eight-layer finetuning job time by dataset. A
submitted job could skip existing fold files, so the distribution is used for
resource accounting rather than architecture-speed claims.}
\label{fig:si-finetune-runtime}
\end{figure*}

\begin{table*}[t]
\centering
\caption{Observed downstream Slurm accounting. A job may skip existing fold files, so these values describe execution history rather than a controlled architecture benchmark.}
\label{tab:si-finetune-runtime}
\resizebox{\textwidth}{!}{%
\begin{tabular}{lrrrrr}
\toprule
Result model & Successful jobs & Other jobs & Successful GPU h & Other GPU h & Median job h \\
\midrule
p3\_base\_vanilla\_n8aligned & 11 & 4 & 53.6 & 9.7 & 4.47 \\
p3\_tpsa\_mofe\_last\_ffn\_n8aligned & 8 & 0 & 71.5 & 0.0 & 1.15 \\
p3\_logp\_mofe\_last\_ffn\_n8aligned & 8 & 0 & 69.1 & 0.0 & 1.18 \\
p5\_tpsa\_mofe\_last\_selected\_kan\_n8aligned & 8 & 1 & 80.4 & 5.2 & 1.21 \\
p5\_logp\_mofe\_last\_selected\_kan\_n8aligned & 8 & 1 & 60.6 & 10.8 & 1.30 \\
p3\_tpsa\_mofe\_last\_ffn\_conloss\_n8aligned & 8 & 0 & 57.7 & 0.0 & 1.21 \\
p3\_logp\_mofe\_last\_ffn\_conloss\_n8aligned & 8 & 1 & 48.1 & 20.7 & 1.17 \\
p5\_tpsa\_mofe\_last\_selected\_kan\_conloss\_n8aligned & 8 & 0 & 71.2 & 0.0 & 1.16 \\
p5\_logp\_mofe\_last\_selected\_kan\_conloss\_n8aligned & 8 & 1 & 38.0 & 24.6 & 1.37 \\
p3\_tpsa\_mofe\_last\_all\_kan\_n8aligned & 10 & 2 & 61.2 & 33.2 & 3.66 \\
p3\_logp\_mofe\_last\_all\_kan\_n8aligned & 8 & 0 & 76.7 & 0.0 & 1.44 \\
\bottomrule
\end{tabular}}
\end{table*}

\section{Metric Definitions and Result Audit}

Table~\ref{tab:si-metric-definitions} records the input, aggregation level,
fixed parameters, and code source for every reported metric. The legacy
eight-layer snapshot contains 440 of 440 planned fold CSV files. The N4 depth reference contains
40 of 40 fold CSV files. The machine-readable result health audit records
file counts, fold and seed mappings, and parsing checks.

\begin{landscape}
\begingroup
\scriptsize
\setlength{\tabcolsep}{3pt}
\renewcommand{\arraystretch}{0.96}
\begin{longtable}{
>{\raggedright\arraybackslash}p{0.11\linewidth}
>{\raggedright\arraybackslash}p{0.27\linewidth}
>{\raggedright\arraybackslash}p{0.14\linewidth}
>{\raggedright\arraybackslash}p{0.17\linewidth}
>{\raggedright\arraybackslash}p{0.15\linewidth}}
\caption{Definitions of evaluation metrics and derived quantities used in the main text and Supporting Information.}
\label{tab:si-metric-definitions}\\
\toprule
Metric & Definition & Aggregation & Parameters & Implementation source \\
\midrule
\endfirsthead
\toprule
Metric & Definition & Aggregation & Parameters & Implementation source \\
\midrule
\endhead
endpoint ROC-AUC & Mask actual == -1 and compute ROC-AUC only when the remaining endpoint contains both classes. & one dataset, fold, and valid endpoint & binary labels and continuous prediction scores & {\ttfamily\tiny analysis\_\allowbreak{}v2/\allowbreak{}build\_\allowbreak{}analysis\_\allowbreak{}v2.py,\allowbreak{} compute\_\allowbreak{}auc} \\
fold ROC-AUC & Arithmetic mean across valid endpoints in one fold. & one dataset and fold & valid endpoints receive equal weight & {\ttfamily\tiny analysis\_\allowbreak{}v2/\allowbreak{}build\_\allowbreak{}analysis\_\allowbreak{}v2.py,\allowbreak{} compute\_\allowbreak{}auc} \\
five-fold mean and sample SD & Arithmetic mean and sample standard deviation with denominator n - 1. & one model and dataset & seeds 42 to 46 and complete 5 of 5 folds & {\ttfamily\tiny analysis\_\allowbreak{}v2/\allowbreak{}build\_\allowbreak{}analysis\_\allowbreak{}v2.py,\allowbreak{} complete\_\allowbreak{}summary} \\
eight-dataset macro ROC-AUC & Equal-weight arithmetic mean of dataset means. & one model across eight classification datasets & no sample-count or endpoint-count weighting & {\ttfamily\tiny analysis\_\allowbreak{}v2/\allowbreak{}build\_\allowbreak{}analysis\_\allowbreak{}v2.py,\allowbreak{} model\_\allowbreak{}macro\_\allowbreak{}summary} \\
paired dataset delta and win & Subtract right from left within fold and average five fold deltas within a dataset. A positive dataset mean is one win. & matched folds followed by equal-weight datasets & identical scaffold seeds for both models & {\ttfamily\tiny analysis\_\allowbreak{}v2/\allowbreak{}build\_\allowbreak{}analysis\_\allowbreak{}v2.py,\allowbreak{} paired\_\allowbreak{}delta\_\allowbreak{}rows} \\
hierarchical paired bootstrap interval & Resample datasets with replacement, then folds with replacement within each sampled dataset. Average folds within dataset and then average datasets. & 10,000 paired hierarchical bootstrap replicates & seed 42 and percentile interval at 2.5 and 97.5 & {\ttfamily\tiny analysis\_\allowbreak{}v2/\allowbreak{}build\_\allowbreak{}analysis\_\allowbreak{}v2.py,\allowbreak{} hierarchical\_\allowbreak{}bootstrap} \\
Wilcoxon signed-rank test & Two-sided Wilcoxon signed-rank test. & one planned comparison & zero\_method=wilcox and Holm adjustment & {\ttfamily\tiny analysis\_\allowbreak{}v2/\allowbreak{}build\_\allowbreak{}analysis\_\allowbreak{}v2.py,\allowbreak{} wilcoxon\_\allowbreak{}holm} \\
conditional interaction & Full - (I1+I2) - (I1+I3) + I1. & matched folds followed by equal-weight datasets & reported separately for TPSA and MolLogP & {\ttfamily\tiny analysis\_\allowbreak{}v2/\allowbreak{}build\_\allowbreak{}analysis\_\allowbreak{}v2.py,\allowbreak{} build\_\allowbreak{}interaction\_\allowbreak{}rows} \\
descriptor-region occupancy & Assign fixed left-closed, right-open intervals with np.digitize and report count and fraction. & one descriptor region & eight fixed regions per descriptor & {\ttfamily\tiny analysis\_\allowbreak{}v2/\allowbreak{}build\_\allowbreak{}analysis\_\allowbreak{}v2.py,\allowbreak{} descriptor\_\allowbreak{}profiles} \\
descriptor-region chemical profile & Region mean and sample SD for each profile variable. & one region and profile variable & regions with fewer than 10 molecules are masked in heatmaps & {\ttfamily\tiny analysis\_\allowbreak{}v2/\allowbreak{}build\_\allowbreak{}analysis\_\allowbreak{}v2.py,\allowbreak{} descriptor\_\allowbreak{}profiles} \\
active pretraining epoch time & Exclude the first epoch of each Slurm segment and values above the segment-specific Q3 + 3 IQR upper outer fence. & median and quartiles within model & queue and inter-job gaps are outside the epoch timer & {\ttfamily\tiny analysis\_\allowbreak{}v2/\allowbreak{}build\_\allowbreak{}analysis\_\allowbreak{}v2.py,\allowbreak{} active\_\allowbreak{}epoch\_\allowbreak{}timing\_\allowbreak{}audit} \\
stored and active model cost & Stored parameters include all branches. Expected active branch cost is occupancy weighted for selected KAN configurations. & one architecture configuration & one MAC equals two FLOPs. KAN excludes spline-basis construction. Routed controls exclude softmax and dispatch overhead & {\ttfamily\tiny analysis\_\allowbreak{}v2/\allowbreak{}build\_\allowbreak{}analysis\_\allowbreak{}v2.py,\allowbreak{} architecture\_\allowbreak{}complexity} \\
pooled molecular embedding & Mean of the first n\_i valid node rows. & one molecule, model, dataset, and fold & n\_i equals the RDKit atom count & {\ttfamily\tiny analysis\_\allowbreak{}v2/\allowbreak{}analyze\_\allowbreak{}embeddings\_\allowbreak{}v2.py,\allowbreak{} mean\_\allowbreak{}pool} \\
embedding reduction and t-SNE & Model-specific PCA to at most 50 dimensions, L2 normalization, then two-dimensional t-SNE for visualization. & one dataset, partition, model, and fold & fold 1, perplexity 30, PCA initialization, learning\_rate=auto, 2000 iterations, and seed 42 & {\ttfamily\tiny analysis\_\allowbreak{}v2/\allowbreak{}analyze\_\allowbreak{}embeddings\_\allowbreak{}v2.py} \\
silhouette score & Mean Euclidean silhouette coefficient. & one model, dataset, partition, and fold & regions with fewer than two samples are excluded & {\ttfamily\tiny analysis\_\allowbreak{}v2/\allowbreak{}analyze\_\allowbreak{}embeddings\_\allowbreak{}v2.py,\allowbreak{} separation\_\allowbreak{}metrics} \\
Davies-Bouldin index & Mean worst-case ratio of within-region dispersion to centroid separation. & one model, dataset, partition, and fold & lower values indicate stronger separation & {\ttfamily\tiny analysis\_\allowbreak{}v2/\allowbreak{}analyze\_\allowbreak{}embeddings\_\allowbreak{}v2.py,\allowbreak{} separation\_\allowbreak{}metrics} \\
inter-to-intra distance ratio & Mean distance between regions divided by mean non-diagonal distance within regions. & one model, dataset, partition, and fold & regions with fewer than two samples are excluded & {\ttfamily\tiny analysis\_\allowbreak{}v2/\allowbreak{}analyze\_\allowbreak{}embeddings\_\allowbreak{}v2.py,\allowbreak{} separation\_\allowbreak{}metrics} \\
\bottomrule
\end{longtable}
\endgroup
\end{landscape}

\section{Embedding Visualizations across Downstream Datasets}

The complete PCA-space metrics are reported in
Table~\ref{tab:si-embedding-all}. The t-SNE figures use fold 1 only and do not
merge coordinate systems across folds. HIV and MUV use deterministic
descriptor-region-stratified samples of at most 2000 common molecules.

\begin{figure*}[t]
\centering
\includegraphics[width=\textwidth]{figS13_embedding_metrics_all_datasets.pdf}
\caption{Five-fold descriptor-region separation metrics across the eight
downstream datasets. A cell is reported only when all five fold embeddings
are available.}
\label{fig:si-embedding-metrics}
\end{figure*}

\begin{figure*}[p]
\centering
\includegraphics[width=\textwidth]{figS14_tsne_bbbp.pdf}
\caption{BBBP fold-1 t-SNE comparison. All panels within a descriptor row use
the same common SMILES set.}
\label{fig:si-tsne-bbbp}
\end{figure*}

\begin{figure*}[p]
\centering
\includegraphics[width=\textwidth]{figS15_tsne_tox21.pdf}
\caption{Tox21 fold-1 t-SNE comparison including the local N8 Base.}
\label{fig:si-tsne-tox21}
\end{figure*}

\begin{figure*}[p]
\centering
\includegraphics[width=\textwidth]{figS16_tsne_toxcast.pdf}
\caption{ToxCast fold-1 t-SNE comparison.}
\label{fig:si-tsne-toxcast}
\end{figure*}

\begin{figure*}[p]
\centering
\includegraphics[width=\textwidth]{figS17_tsne_sider.pdf}
\caption{SIDER fold-1 t-SNE comparison.}
\label{fig:si-tsne-sider}
\end{figure*}

\begin{figure*}[p]
\centering
\includegraphics[width=\textwidth]{figS18_tsne_clintox.pdf}
\caption{ClinTox fold-1 t-SNE comparison.}
\label{fig:si-tsne-clintox}
\end{figure*}

\begin{figure*}[p]
\centering
\includegraphics[width=\textwidth]{figS19_tsne_bace.pdf}
\caption{BACE fold-1 t-SNE comparison.}
\label{fig:si-tsne-bace}
\end{figure*}

\begin{figure*}[p]
\centering
\includegraphics[width=\textwidth]{figS20_tsne_hiv.pdf}
\caption{HIV fold-1 t-SNE comparison using a deterministic
descriptor-region-stratified sample of the common molecules.}
\label{fig:si-tsne-hiv}
\end{figure*}

\begin{figure*}[p]
\centering
\includegraphics[width=\textwidth]{figS21_tsne_muv.pdf}
\caption{MUV fold-1 t-SNE comparison using a deterministic
descriptor-region-stratified sample of the common molecules.}
\label{fig:si-tsne-muv}
\end{figure*}

\section{Machine-Readable Analysis Artifacts}

The canonical data package includes fold-level AUC values, complete-cell
summaries, paired contrasts, descriptor-region profiles, KAN screening
results, embedding coordinates and metrics, active-time audits, Slurm
accounting, metric definitions, result health checks, and a figure-table
provenance manifest. Each record retains its source result or log path.



\end{document}
